% =======================================================================
%  Branch-Local Hilbert-Polya:
%  A Boundary Theorem from Riemann, Selberg, and Prime-Hub
%
%  Autor: Lukas Geiger, Unabhängiger Forscher, Bernau
%  Quellencheck-Revision v2 - 2026-05-14
%  Status: quellenkorrigierter bilingualer Kandidat
% =======================================================================
\documentclass[11pt]{article}

\usepackage[letterpaper,top=2.4cm,bottom=2.6cm,left=2.2cm,right=2.2cm]{geometry}
\usepackage[utf8]{inputenc}
\usepackage{cmap}
\usepackage[T1]{fontenc}
\usepackage{lmodern}
\usepackage{microtype}
\setlength{\emergencystretch}{2em}
\usepackage[ngerman]{babel}
\usepackage{amsmath,amssymb,amsfonts,mathtools,bm}
\usepackage{graphicx,float,booktabs,array,multirow}
\usepackage{xcolor}
\usepackage[unicode=true,colorlinks=true,linkcolor=blue!60!black,citecolor=blue!60!black,urlcolor=blue!60!black]{hyperref}
\hypersetup{
  pdftitle={Der Zeta-Zoo: Die mathematische Seite der funktionalen Stabilitätstheorie},
  pdfauthor={Lukas Geiger},
  pdfsubject={Branch-lokale Hilbert-Polya-Klassifikation und Quellencheck-Revision},
  pdfkeywords={Funktionale Stabilitätstheorie, Zeta-Zoo, Hilbert-Polya, Weil-Form, SGE, UCU}
}

\usepackage{amsthm}
\newtheorem{theorem}{Theorem}[section]
\newtheorem{proposition}[theorem]{Proposition}
\newtheorem{corollary}[theorem]{Korollar}
\newtheorem{lemma}[theorem]{Lemma}
\newtheorem{conjecture}[theorem]{Vermutung}
\theoremstyle{definition}
\newtheorem{definition}[theorem]{Definition}
\newtheorem{example}[theorem]{Beispiel}
\theoremstyle{remark}
\newtheorem{remark}[theorem]{Bemerkung}

\newcommand{\HBL}{\mathrm{HP\text{-}BL}}
\newcommand{\PSL}{\mathrm{PSL}}
\newcommand{\Spec}{\operatorname{Spec}}
\renewcommand{\Re}{\operatorname{Re}}
\newcommand{\halfline}{\tfrac{1}{2}}
\newcommand{\zer}[1]{\zeta'/\zeta(#1)}
\DeclareMathOperator{\Weil}{Weil}

\title{Der Zeta-Zoo\\
  \large --- Die mathematische Seite der\\
  funktionalen Stabilitätstheorie\\[0.3em]
  \normalsize Branch-Lokales Hilbert-P\'olya via die Trinität\\
  von UCU, SGE und Weil-Form-Transversalität}
\author{Lukas Geiger\\
  \small Unabhängiger Forscher, Bernau\\
  \small \href{https://orcid.org/0009-0005-7296-1534}{ORCID: 0009-0005-7296-1534}}
\date{\today}

\begin{document}
\maketitle

\begin{abstract}
\noindent
Die funktionale Stabilitätstheorie (FST) ist ein Programm, das ein
einziges strukturelles Muster --- funktionale Positivität unter einer
Eichbedingung (Gauge Constraint) --- als gemeinsames Substrat offener Probleme in
der Zahlentheorie, der mathematischen Physik und der Kosmologie identifiziert \cite{RFEPFramework}.
\textbf{Dieses Paper ist die mathematische Seite der FST.} Wir formalisieren
die Trinität (Dreifaltigkeit) von Metaprinzipien, die den Zeta-Typ-Zweig der
Theorie steuern: das Lemma der \emph{universellen Konvexitätseindeutigkeit} (UCU),
das die strikte Konvexität des maßgeblichen Funktionals in den
eindeutigen Sättigungszustand umwandelt; die Vermutung der
\emph{Halbgruppen-Gruppen-Äquivalenz} (SGE), die die Hilbert-P\'olya-Eigenschaft
als Kandidatin für eine binäre algebraische Invariante der Transfer-Provenienz
behandelt; und die zweigübergreifende \emph{Weilsche quadratische Form}, die eine
Lücken-Positivität (gap positivity) auf beiden Seiten der SGE-Dichotomie
liefert. Innerhalb dieser Trinität reorganisiert sich das klassische Bild
sauber: die Nichtexistenz-Resultate NE-A und NE-B des Riemann-v2.1-Programms
\cite{Geiger2026RHII} schließen die algebraischen Routen für $\zeta(s)$ aus;
Selbergs Laplace-Beltrami-Operator realisiert die Hilbert-P\'olya-Struktur
explizit für $Z_\Gamma(s)$; der spektrale Zookeeper \cite{Zookeeper2026}
schließt den Riemannschen $C^{(2)}$/MS2-Schritt im CCM/Fourier-Zweig ab; der
konjekturale Prime-Hub-Graph bleibt offen als eine explizite OP5-Konstruktion
unter vier dokumentierten strukturellen Hindernissen, einschließlich des
Determinanten-Orientierungshindernisses zwischen $1/\zeta(s)$ und
Eigenwert-eins-Ereignissen.
Wir klassifizieren den \emph{Zeta-Zoo} entlang der SGE-Dichotomie und ergänzen
endliche Kalibrierungschecks an fünf Testfamilien (Riemann, Selberg,
Prime-Hub, Dedekind $\mathbb{Q}(\sqrt{-5})$, Ihara-Petersen). Diese Checks
stützen die Taxonomie, beweisen aber weder SGE noch einen globalen
Hilbert--P\'olya-Transfer. Außerdem leiten wir
Konsequenzen für Dirichletsche $L$-Funktionen ab \cite{PaleyWienerDirichlet, DirichletAtlas}.
Wir schließen mit einer kryptographischen Lesart der Trinität als der algebraischen
Architektur des Bitcoin-Blockchain-Protokolls und mit dem narrativen Kreis
vom Kooperativen Renormierungsmodell (CRM) der Dunklen Energie durch die FST
hin zur abstrakten Trinität und zurück via hyperbolischer Geometrie. Die physikalische
Instanziierung des FST-Programms über das renormierte Prinzip der freien Energie (RFEP)
wird im Begleitpaper \cite{RFEPFramework} entwickelt.
\end{abstract}

\clearpage
\tableofcontents
\clearpage

% ============================================================
\section{Einleitung}
\label{sec:intro}

\subsection{Die Hilbert-P\'olya-Vermutung und ihre implizite Zweig-Wahl}
\label{sec:hp-implicit-branch}

Die Hilbert-P\'olya-Vermutung fragt danach, ob die nicht-trivialen Nullstellen der
Riemannschen Zetafunktion den Eigenwerten eines selbstadjungierten Operators
entsprechen. Ihr ältester dokumentierter Ursprung ist ein Brief von G.~P\'olya
an A.~Odlyzko aus dem Jahr 1982 \cite{Polya1982letter}, der eine Unterhaltung
mit E.~Landau in Göttingen um 1912--1914 wiedergibt; die erste
veröffentlichte Erwähnung erscheint erst in Montgomerys Paper zur Paar-Korrelation
von 1973 \cite{Montgomery1973}; Standardüberblicke und Millennium-Problem-
Darstellungen bieten \cite{Bombieri2000,Conrey2003}; und die Vermutung wurde
seitdem wiederholt aufgegriffen
\cite{Berry1985, BerryKeating1999, KeatingSnaith2000, Connes1999}. Sie
hat einen Großteil der analytischen Zahlentheorie des letzten Jahrhunderts
angetrieben. Dennoch wurde, wie Connes' Übersicht von 2026 \cite{Connes2026}
betont, ein solcher Operator für $\zeta(s)$ bisher nicht explizit konstruiert.

Das Programm trifft implizit eine Zweig-Wahl (branch choice). Für die Selbergsche Zetafunktion
$Z_\Gamma(s)$ einer kompakten hyperbolischen Fläche
$\Gamma\backslash\mathbb{H}$ realisiert der Laplace-Beltrami-Operator
$\Delta_{\Gamma\backslash\mathbb{H}}$ die Hilbert-P\'olya-Struktur
explizit: seine Eigenwerte $\lambda_n = 1/4 + r_n^2$ liefern
$r_n = \operatorname{Im}(\rho_n)$ für die Selberg-Nullstellen $\rho_n = 1/2 + ir_n$,
und er ist mit der kanonischen Spurformel-Transferalgebra kompatibel, die durch
den geodätischen Isometriefluss erzeugt wird \cite{Selberg1956}. Für die
Riemannsche Zetafunktion hingegen
schließen die Resultate NE-A und NE-B aus dem jüngsten Even-Dominance-Beweisprogramm
\cite{Geiger2026RHII} die beiden natürlichsten algebraischen Kandidaten aus:
der Primverschiebung-Fouriermultiplikator
$M(\xi) = -2\,\Re\zer{\halfline + i\xi}$ scheitert daran, auf einer Menge von positivem
Maß absolut total-positiv zu sein (NE-A), und keine nicht-skalare symmetrische Matrix
kommutiert simultan mit allen Primverschiebungsmatrizen bei einer
Galerkin-Trunkierung von $N\leq 15$ (NE-B, wobei die Erweiterung auf den vollen Raum konjektural
ist \cite{Geiger2026RHII}). Und für den konjekturalen Prime-Hub-Graphen (Offenes Problem 5 in
\cite{Geiger2026RHIII}) halten vier Hindernisse --- Inkompatibilität des Weylschen
Gesetzes mit der logarithmischen Dichte $N(T)\sim (T/2\pi)\log(T/2\pi)$, Spurklasse-Regularität
eines Graphentransferoperators, Kompatibilität mit NE-A und die
Determinantenorientierung zwischen $1/\zeta(s)$ und gewöhnlichen
Fredholm-Nullen --- eine explizite
Prime-Hub/Hilbert-P\'olya-Graphenkonstruktion offen. Nach dem spektralen Zookeeper \cite{Zookeeper2026}
ist diese Offenheit nicht länger der Status des Riemannschen $C^{(2)}$/MS2-Schritts
selbst; es ist der Status der direkten OP5-Grapheninterpretation.

\subsection{Die Trinität der Metaprinzipien}
\label{sec:trinity}

Diese drei kanonischen Fakten deuten gemeinsam darauf hin, dass die Hilbert-P\'olya-Eigenschaft
kein universelles Merkmal von Zeta-Typ-Familien ist, sondern eine \emph{lokale}
Invariante, die in manchen Unterfamilien präsent und in anderen strukturell
abwesend ist. Das zentrale organisierende Prinzip dieses Papers ist,
dass diese Branch-Lokalität durch eine Trinität von Metaprinzipien gesteuert wird,
die jedes Programm vom RH-Typ adressieren muss:

\begin{description}
\item[(UCU)] \emph{Universelle Konvexitätseindeutigkeit (Universal Convexity Uniqueness).} Die zulässigen
  Zustände des Programms sind die Minimierer eines strikt konvexen,
  koerziven Funktionals. Strikte Konvexität garantiert einen \emph{eindeutigen}
  Sättigungszustand, den die spektrale Maschinerie der Weilschen quadratischen
  Form dann realisiert. Dieses Prinzip organisiert die Eindeutigkeitsargumente
  über vier ansonsten unverbundene Programme (Turbulenz~K41, FST Dunkle-Energie-Potenzial,
  Riemannsche Even-Dominance, Yang--Mills Gibbs-Maß) unter einem einzigen Dach.

\item[(SGE)] \emph{Halbgruppen-Gruppen-Äquivalenz (Semigroup-Group Equivalence).} Ob ein natürlicher
  selbstadjungierter Operator --- der gesuchte Hilbert-P\'olya-Hamiltonian
  $H_X$ --- existiert, wird konjektural durch eine einzige algebraische
  Provenienz-Invariante gesteuert: ob die Transferstruktur lokalkompakte
  Gruppen-Provenienz oder nur halbgruppenartige Provenienz besitzt. In den
  realisierten geometrischen Fällen ist dies die Bedingung, unter der eine
  Lie-Struktur und damit ein Casimir-Operator hervortritt.

\item[(WF)] \emph{Weil-Form Branch-Transversalität.} Die Weilsche quadratische Form $Q_X$ ist
  für jede Zeta-Typ-Familie mit einem Euler-Produkt, einem Gamma-Faktor und
  einer Funktionalgleichung wohldefiniert, unabhängig vom $\HBL$-Status. Sie
  liefert eine einheitliche Architektur, die Lücken-Positivität (gap positivity) in allen
  algebraischen Regimes erzeugt. Im Riemannschen Fall wird das
  $Q_\zeta$-basierte Companion-Programm der Direkten Frontier-Dominanz
  \cite{Geiger2026RHII} derzeit als bedingte Reduktions- und
  Diagnoseroute zur Riemannschen Vermutung geführt, ohne jeglichen
  Hilbert-P\'olya-Input.
\end{description}

Kurz gesagt:

\begin{quote}
\emph{UCU selektiert den eindeutigen Sättigungszustand; SGE entscheidet, ob
ein natürlicher Hilbert-P\'olya-Operator existiert; die Weil-Form-Transversalität
baut die Brücke über die SGE-Dichotomie hinweg.}
\end{quote}

Der Rest dieses Papers formalisiert diese Trinität, testet sie an
fünf Zeta-Typ-Familien und schließt mit einer spekulativen Lesart, in der
die drei Prinzipien die algebraische Architektur des Bitcoin-Blockchain-Protokolls
widerspiegeln (Abschnitt~\ref{sec:blockchain}).

\subsection{Ein Exkursionsführer für den Zeta-Zoo}
\label{sec:zoo-framing}

In diesem Paper verwenden wir den \emph{Zeta-Zoo} als pädagogischen
Rahmen. Jede Zeta-Typ-Familie $X$ ist ein Gehege im Zoo; die Gehegearten sind die drei möglichen
HP-Status-Werte (\text{YES}, \text{NO}, \text{OPEN}); das zooweite Futter ist die
Weilsche quadratische Form $Q_X$, die jedes Gehege unabhängig von der Art annimmt;
der Zooführer ist die Trinität UCU + SGE + Weil-Form-Transversalität und die
Vier-Klassen-Rigoroshierarchie (Abschnitt~\ref{sec:rigor}). Modellgehege (Riemann,
Selberg, Prime-Hub) sind die Hauptstationen der Tour. Neuzugänge
(Abschnitt~\ref{sec:sge-dedekind-test}: Dedekind $\mathbb{Q}(\sqrt{-5})$;
Abschnitt~\ref{sec:sge-ihara-test}: Ihara-Petersen) dienen als erste nicht-kanonische Zoo-Exponate.
Spezies, die auch das benachbarte FST-physikalische Biotop
\cite{RFEPFramework} bewohnen --- \emph{Brückenspezies} wie die Riemannsche Zetafunktion selbst,
die BSD L-Funktion, das Hodge-Motiv und die boolesche Algebra von P vs NP ---
sind an den entsprechenden Einträgen querverwiesen.

Die Zoo-Metapher ist eine Lesehilfe, kein theorematisches Gerüst. Alle
untenstehenden mathematischen Aussagen sind technisch formuliert; die
Metapher wird lediglich verwendet, um den Leser durch die Vielfalt der Zeta-Typ-Familien
zu orientieren, die die Trinität organisiert.

\subsection{Der CRM-Kreis: vom Kosmos zur Algebra und zurück}
\label{sec:crm-circle}

Historisch gesehen war das Kooperative Renormierungsmodell (CRM) der
kosmologischen Dunklen Energie die erste Instanziierung dessen, was später zur
funktionalen Stabilitätstheorie wurde \cite{RFEPFramework}. Seine
Kompositionsregel
\[
  x \oplus y = \frac{x + y}{1 + xy}
\]
ist die Lorentz-Boost-Addition, sodass die zugrundeliegende Geometrie
konstruktionsbedingt hyperbolisch ist. Als das CRM-Stabilitätsargument
generalisiert wurde, entstand die FST; als die FST weiter abstrahiert wurde,
kristallisierte sich die mathematische Trinität von UCU, SGE und Weil-Form-Transversalität
heraus. Der Kreis schließt sich: vom Kosmos zur Physik zum Prinzip zur Algebra,
und die hyperbolische Geometrie, die das CRM konkret realisiert,
taucht auf der YES-Seite der SGE-Dichotomie als der $\PSL_2(\mathbb{R})$-Casimir
der Selberg-Zeta (Abschnitt~\ref{sec:selberg}) wieder auf und, auf beiden Seiten,
als die branch-transversale Weil-Form-Architektur (Abschnitt~\ref{sec:weil-transversal}).
Am Ende ist alles Algebra --- einschließlich der Geometrie.

\subsection{Organisation}
\label{sec:intro-organization}

Abschnitt~\ref{sec:setup} ruft die Definitionen der Hilbert-P\'olya-Struktur und
der Weilschen quadratischen Form im Kontext der Zeta-Familien in Erinnerung.
Abschnitt~\ref{sec:locality} führt die Branch-Lokalität formal ein und bettet sie
in die Trinität ein. Die Abschnitte \ref{sec:riemann}, \ref{sec:selberg} und \ref{sec:primehub}
erinnern an die drei kanonischen Instanzen: Riemann (NE-A und NE-B), Selberg
(Casimir-Operator) und Prime-Hub (Offenes Problem 5).
Abschnitt~\ref{sec:boundary} formuliert und beweist das Randtheorem (boundary theorem).
Abschnitt~\ref{sec:ucu} führt UCU formal ein und zeigt seine
Rolle als die universelle Eindeutigkeitsaussage über vier Programme hinweg auf.
Abschnitt~\ref{sec:weil-transversal} diskutiert Weil-Form-Methoden als das
branch-transversale Komplement. Abschnitt~\ref{sec:sge-conjecture}
formuliert die SGE-Vermutung und testet sie an den Dedekind- und
Ihara-Familien. Abschnitt~\ref{sec:blockchain} schließt mit der
Blockchain-Lesart, die die Essenz der Trinität in ein informatisches Register
überträgt. Abschnitt~\ref{sec:outlook} sammelt die Rigoros-Hierarchie, die Konsequenzen
für Dirichlet und den Ausblick auf angrenzende Projekte.

% ============================================================
\section{Setup}
\label{sec:setup}

\subsection{Zeta-Typ-Familien}
\label{sec:zeta-family}

\begin{definition}[Zeta-Typ-Familie]\label{def:zeta-family}
Eine \emph{Zeta-Typ-Familie} ist ein Tupel $(Z_X, \mathcal{T}_X, \mathcal{E}_X)$
bestehend aus
\begin{enumerate}
\item einer meromorphen Funktion $Z_X(s)$ auf $\mathbb{C}$ mit höchstens
  endlich vielen Polen, nicht-trivialen Nullstellen $\{\rho_k = 1/2 + i\gamma_k^X\}$
  auf oder in der Nähe einer kritischen Gerade $\{\Re(s) = 1/2\}$, die ein Euler-Produkt
  auf einer rechten Halbebene und eine Funktionalgleichung
  $Z_X(s) \leftrightarrow Z_X(1-s)$ bis auf einen Gamma-Faktor erfüllt;
\item einer natürlichen Familie $\mathcal{T}_X$ von beschränkten Operatoren auf einem
  Hilbertraum, die mit der Geometrie/Arithmetik von $X$ assoziiert ist
  ("Transfer-Operatoren");
\item einer expliziten Formel $\mathcal{E}_X$, die Testfunktions-Summen
  über Nullstellen mit Prim-/Geodäten-Summen der Familie in Beziehung setzt.
\end{enumerate}
\end{definition}

Für Euler-Produkte, Funktionalgleichungen und automorphe $L$-Funktionen
verwenden wir die Standardterminologie aus \cite{IwaniecKowalski}.

\begin{example}[Kanonische Instanzen]\label{ex:canonical}
(i) $X = \zeta$: $\mathcal{T}_\zeta = \{\text{Primverschiebungs-Operatoren } D(r)\}$ mit expliziter Formel vom Weil-Typ.\\
(ii) $X = Z_\Gamma$ für $\Gamma$ ko-kompakt Fuchssch: $\mathcal{T}_\Gamma$
ist die kanonische Spurformel-/Wave-Kernel-Algebra, die durch den
geodätischen Fluss induziert wird, mit Selbergs expliziter Formel.\\
(iii) $X = G_\mathrm{prime}$: konjekturaler Transferoperator auf einem
primzahl-indizierten Graphen; Offenes Problem 5 in \cite{Geiger2026RHIII}.
\end{example}

\subsection{Die Hilbert-P\'olya-Struktur}

\begin{definition}[Branch-Lokales Hilbert-P\'olya]\label{def:hp-bl}
Man sagt, eine Zeta-Typ-Familie $(Z_X, \mathcal{T}_X, \mathcal{E}_X)$
\emph{besitzt Hilbert-P\'olya-Struktur} --- geschrieben $\HBL(X) = \mathrm{YES}$ ---
wenn ein selbstadjungierter Operator $H_X : \mathcal{D}(H_X) \to \mathcal{H}_X$ auf
einem Hilbertraum $\mathcal{H}_X$ existiert, sodass:
\begin{enumerate}
\item[(HP1)] $\Spec(H_X) = \{\gamma_k^X : k \in \mathbb{N}\}$, mit Vielfachheiten.
\item[(HP2)] $H_X$ kommutiert mit jedem $T \in \mathcal{T}_X$ im Sinne
  kommutierender Spektralprojektionen.
\item[(HP3)] $H_X$ ist \emph{natürlich}: eindeutig determiniert durch die
  Geometrie oder Arithmetik von $X$, bis auf unitäre Äquivalenz.
\end{enumerate}
Existiert kein solches $H_X$, schreiben wir $\HBL(X) = \mathrm{NO}$;
ist die Existenz offen, $\HBL(X) = \mathrm{OPEN}$.
\end{definition}

\begin{remark}[Informelle Intuition]\label{rem:hp-informal}
Hilberts und P\'olyas ursprünglicher Vorschlag war, dass die Riemannsche
Vermutung aus dem Vorweisen eines natürlichen selbstadjungierten Operators
"in der Natur" folgen würde, dessen Eigenwerte mit den Imaginärteilen der
Zeta-Nullstellen übereinstimmen. Die drei Bedingungen (HP1)--(HP3)
formalisieren diese informelle Anforderung: spektrale Identität, Kommutierung mit den
natürlichen Transfer-Operatoren der Familie, und eine Natürlichkeitsbedingung,
die tautologische Konstruktionen (z.B. den trivialen Multiplikationsoperator
auf $\ell^2$, der durch die Nullstellen selbst indiziert wird) verhindert.
\end{remark}

\subsection{Die Weilsche quadratische Form als branch-transversale Architektur}
\label{sec:weil-qf}

Für jede Zeta-Typ-Familie $(Z_X, \mathcal{T}_X, \mathcal{E}_X)$ mit
vervollständigter Funktionalgleichung und einer expliziten Formel im
üblichen Weil-Sinn induziert die Weilsche explizite Formel
\cite{Weil1952} eine quadratische Form $Q_X$ auf der gewählten
Testfunktionsklasse von der schematischen Gestalt
\begin{equation}
\label{eq:weil-qf-generic}
  Q_X[f] = \sum_k 2\pi |\widehat{f}(\gamma_k^X)|^2
    = W^X_\infty[f] + (\text{Prim-/Geodätenseite}).
\end{equation}
Nachdem die Standardkonventionen der Vervollständigung festgelegt sind
(Gamma-Faktor, Pole, triviale Nullstellen und außergewöhnliche
Spektralterme), ist Positivität von $Q_X$ auf der relevanten
Testfunktionsklasse die Weil-Kriteriumsform der RH-artigen Aussage für
$Z_X$, unabhängig von der Existenz irgendeines Hilbert-P\'olya-Operators.
Dies ist der branch-transversale Inhalt der Weilschen Architektur: er
ist in den Fällen HP-BL-YES, -NO und -OPEN anwendbar, aber nur nachdem
diese lokalen Normalisierungen spezifiziert wurden.

\subsection{Dreikomponenten-Zerlegung von RH-Typ-Aussagen}
\label{sec:three-component}

Eine weitere strukturelle Verfeinerung isoliert im Anschluss an
\cite{MetaGalerkinBridge} drei konzeptionell verschiedene
Schichten (layers) innerhalb jeder RH-Typ-Aussage für eine Zeta-Typ-Familie $X$.
Schreibt man $\mathrm{RH}(X)$ für die Positivität von $Q_X$ auf der
zulässigen Testfunktionsklasse, verwenden wir das folgende
Beweispflichten-Routenledger, nicht eine Theorem-Äquivalenz:
\begin{equation}
\label{eq:three-component}
  \mathsf{RHRoute}(X):\qquad
  \underbrace{\mathrm{GBC}(X)}_{\text{Galerkin-Brücke}}
  \;\longrightarrow\;
  \underbrace{C^{(2)}(X)}_{\text{Eigenvektor-Realisierung}}
  \;\longrightarrow\;
  \underbrace{\mathrm{Limit}_{\mathrm{Hurwitz}}(X)}_{\text{Universeller Limes}}.
\end{equation}
\eqref{eq:three-component} hält fest, welche drei Schichten jedes Programm
vom RH-Typ liefern muss, und kennzeichnet diese als Ziele getrennter
technischer Unterprogramme. In jeder konkreten Familie müssen die Labels mit
eigenen Definitionsbereichs-, Kompaktheits-, Konvergenz- und
Nichtverschwindenshypothesen instanziiert werden; die präzisen Konventionen
werden in den jeweiligen Instanzen der
Abschnitte~\ref{sec:riemann}--\ref{sec:primehub} festgelegt.

Die drei Komponenten lauten wie folgt.

\begin{enumerate}
\item[\textbf{(GBC)}] \emph{Galerkin-Brücken-Komponente.} Der
algebraisch-spektrale Schritt: für jeden Cutoff $\lambda > 1$ konstruiert
man eine endlich-dimensionale Galerkin-Approximation $Q_X^{(N,\lambda)}$
von $Q_X$ auf einem logarithmischen Host-Raum (Kosinus-/Sinusbasis auf
$[-\log\lambda,\log\lambda]$ im arithmetischen Fall), zeigt die
Positivität/Negativität der Approximation über Intervallarithmetik
oder kombinatorische Argumente, und überträgt das Resultat über Mosco-Konvergenz
\cite{MetaGalerkinBridge,ConnesConsaniMoscovici2025} auf den kontinuierlichen Operator.
  Diese Schicht ist der Ort, an dem die v2.1 Even-Dominance- /
  Direct-Frontier-Reduktionsroute \cite{Geiger2026RHII} für
  $X = \zeta$ lebt.

\item[$\boldsymbol{C^{(2)}}$] \emph{Eigenvektor-Realisierungs-Komponente.} Der analytische
Schritt: man zeigt, dass die Galerkin-Eigenvektoren $f_\star^{(N)}$ in einer
ausreichend starken Topologie gegen eine echte Eigenfunktion $f_\star$ des kontinuierlichen
Operators konvergieren. Für den Riemannschen Fall ist dies genau die Paley--Wiener-Kontrolle,
die in Connes' Programm benötigt wird
\cite{Connes2026,ConnesConsaniMoscovici2025} und in \cite{Geiger2026RHII}
als der tiefste verbleibende analytische Baustein vor dem Zookeeper-Abschluss gekennzeichnet
wurde. Im aktuellen CORE-Status schließt \cite{Zookeeper2026} diesen
$C^{(2)}$/MS2-Schritt im CCM/Fourier-Modell ab und liefert die unten
verwendete Mikrocluster-Normalform. Für Selberg ist dieser Schritt automatisch (kompakte Resolvente
des Laplace-Operators auf $\Gamma\backslash\mathbb{H}$).

\item[\textbf{(Hurwitz)}] \emph{Universelle Limes-Komponente.} Der
letzte funktionale Schritt: Der Satz von Hurwitz über lokal gleichmäßige Grenzwerte
nicht verschwindender holomorpher Funktionen hebt die Familie von
Level-Aussagen in den Limes auf der kritischen Gerade. Diese
Komponente ist unabhängig von $X$: es ist dasselbe analytische Lemma,
das in jedes $\lambda$-Cutoff-RH-Argument einfließt, weshalb wir es
separat als "universell" auflisten.
\end{enumerate}

\begin{remark}[Wo sich jeder FST-Zweig spezialisiert]
\label{rem:three-component-specialisation}
Die Dreikomponenten-Zerlegung macht sichtbar, \emph{wo} jeder FST-Zweig seinen
technischen Fokus setzt:
\begin{itemize}
\item Das \textbf{v2.1 Even Dominance Programm} \cite{Geiger2026RHII}
wird derzeit als bedingte Reduktions- und Diagnoseroute für
$X = \zeta$ geführt: endliche CAP-Rechnungen und die GBC-Schicht sind
Evidenz- und Normalisierungsdaten, während die öffentliche RH-Folgerung
von einem Odd-Sector-Lower-Bound / AVG-artigen Abschluss und der
zugehörigen $C^{(2)}(\zeta)$-Brücke abhängt.
\item Das \textbf{Connes-Programm} \cite{Connes2026,ConnesConsaniMoscovici2025}
ist die natürliche Heimat des $C^{(2)}$-Schritts, der semilokal als eine analytische
Bedingung an prolate-sphäroidische Eigenfunktions-Asymptotiken formuliert ist;
im Zookeeper Paper ist dieser Schritt im CCM/Fourier-Zweig durch das Drei-Lemma-Mikrocluster-Argument
abgeschlossen.
\item Das \textbf{Hindernis-Klassifikator-Programm} (Obstruction-Classifier)
\cite{GeigerObstructionClassifiers} diagnostiziert, an welchem
(GBC, $C^{(2)}$)-Paar ein gegebenes Problem schwierig ist: leichtes GBC +
Zookeeper-geschlossenes $C^{(2)}$ (Riemann), twisted-transfer
$C^{(2)}$ (Dirichlet), beides leicht (Selberg, automorphe
$L$ auf $GL_n$), offenes GBC (Ihara generisch, Yang--Lee, $p$-adische
$L$).
\item Die \textbf{Hurwitz}-Schicht ist über alle diese Zweige invariant; wenn wir
"den Limes nehmen" zur kontinuierlichen kritischen Gerade, verrichtet dasselbe
Lemma die Arbeit.
\end{itemize}
Insbesondere wirkt die Trinität von Abschnitt~\ref{sec:ucu} --- UCU + SGE +
Weil-Form-Transversalität --- auf die GBC-Schicht (UCU liefert die
Eindeutigkeit des Galerkin-Grundzustands, SGE klassifiziert die zulässigen
Transfer-Algebren, die Weil-Form-Transversalität ist die Architektur,
die die Galerkin-Approximation trägt). $C^{(2)}$ und Hurwitz sind zu dieser Trinität
transversal; sie kommen ins Spiel, \emph{nachdem} die Trinität den
Grundzustand festgelegt hat.
\end{remark}

\begin{remark}[Die Randtheorem-Sichtweise]
\label{rem:boundary-three-component}
Theorem~\ref{thm:boundary} (HP-BL) klassifiziert Zeta-Typ-Familien
nach ihrem HP-Status. Die Dreikomponenten-Zerlegung verfeinert dies,
indem sie uns mitteilt, \emph{welche Komponente die Klassifizierung determiniert}:
\begin{itemize}
\item HP-BL $=$ YES (Selberg, Ihara Petersen, automorph):
$C^{(2)}$ ist automatisch (kompakte Resolvente); der HP-Operator ist
der natürliche Laplace-Operator.
\item HP-BL $=$ NO (Riemann, Dedekind $\mathbb{Q}(\sqrt{-5})$, Dirichlet):
Ein Fehlschlag/Nichtkonstruktion von HP-BL ist nicht identisch mit einem
offenen $C^{(2)}$-Status. Für Riemann schließt das Zookeeper Paper
$C^{(2)}$/MS2 im CCM/Fourier-Zweig ab; für Dirichlet und Dedekind muss der
entsprechende getwistete (verdrehte) Transfer noch getestet werden. Die Weil-Form-Route
bleibt branch-transversal.
\item HP-BL $=$ OPEN (Prime-Hub): die explizite OP5 Graphen-Konstruktion
und ihre Randoperator-Interpretation bleiben offen; nach Zookeeper wird dies
nicht länger als offener Riemann-Blocker für $C^{(2)}$ gelesen.
\end{itemize}
\end{remark}

% ============================================================
\section{Branch-Lokalität als Konzept}
\label{sec:locality}

Branch-Lokalität bezieht sich auf die Beobachtung, dass eine strukturelle
Eigenschaft von Zeta-Typ-Familien für die eine Familie gelten und für eine eng
verwandte fehlschlagen kann. Bevor wir das Randtheorem formulieren, machen
wir das Konzept präzise.

\subsection{Drei Modi von \texorpdfstring{$\HBL$}{HP-BL}}

\begin{definition}[HP-BL Status-Trichotomie]\label{def:hpbl-status}
Für eine Zeta-Typ-Familie $X$ im Sinne von Definition~\ref{def:zeta-family}
sagen wir, dass
\begin{itemize}
\item $\HBL(X) = \mathrm{YES}$, wenn ein Operator $H_X$ existiert, der
  (HP1)--(HP3) aus Definition~\ref{def:hp-bl} erfüllt;
\item $\HBL(X) = \mathrm{NO}$, wenn man beweisen kann oder unter einer
  spezifischen Vermutung etabliert hat, dass kein solches $H_X$ existiert;
\item $\HBL(X) = \mathrm{OPEN}$, wenn kein Kandidat vorgezeigt wurde
  \emph{und} kein Ausschluss bewiesen wurde, möglicherweise vorbehaltlich
  identifizierter Hindernisse.
\end{itemize}
\end{definition}

Der Status $\mathrm{NO}$ kann \emph{unbedingt} (unconditional) oder
\emph{bedingt} (conditional) sein. In der vorliegenden Arbeit wird der Riemannsche
Fall bedingt nach Vermutung~\ref{conj:NE-B-full} (der Ausweitung von NE-B
von endlichdimensionalen Galerkin-Trunkierungen auf den vollen Raum) $\mathrm{NO}$
sein. Die Vermutung ist konjektural, aber strukturell präzise: eine explizite
offene Aussage, deren Auflösung die NO-Klassifikation entweder bestätigen
oder widerlegen würde.

\subsection{Natürlichkeit (HP3): der operationale Test}

Die Natürlichkeitsklausel (HP3) von Definition~\ref{def:hp-bl} ist die
subtile Komponente. Ohne sie existiert stets ein trivialer "HP-Operator": auf
dem Hilbertraum $\ell^2(\{\gamma_k\})$ kann man $H_X$ einfach als
Multiplikation mit $\gamma_k$ definieren. Ein derartiges ad-hoc $H_X$ erfüllt
(HP1) tautologisch, erklärt aber gar nichts. Wir verlangen daher,
dass $H_X$ im folgenden operationalen Sinn \emph{natürlich} ist.

\begin{definition}[Natürlicher HP-Operator]\label{def:natural-hp}
Ein selbstadjungierter Operator $H_X$ auf $\mathcal{H}_X$ wird
\emph{natürlich} für die Zeta-Typ-Familie $X$ genannt, wenn:
\begin{enumerate}
\item[(N1)] es eine natürliche Symmetriegruppe $G_X$ gibt, die auf $X$ operiert
  (über dessen Geometrie, Arithmetik oder Automorphismenstruktur), oder eine
  vergleichbar kanonische Operator-Provenienzstruktur wie eine Hecke-Algebra,
  ein Gruppoid, eine adelische Darstellung oder ein Spurformel-Paket, unter der
  $H_X$ invariant ist;
\item[(N2)] der Hilbertraum $\mathcal{H}_X$ und die Einbettung
  $\mathcal{T}_X \hookrightarrow \mathcal{B}(\mathcal{H}_X)$ $G_X$-äquivariant sind;
\item[(N3)] $H_X$ durch (HP1)--(HP2) und (N1)--(N2) bis auf unitäre Äquivalenz
  kommutierend mit der $G_X$-Operation eindeutig determiniert ist.
\end{enumerate}
\end{definition}

In Selbergs Fall ist die Symmetriegruppe $\Gamma\backslash\PSL_2(\mathbb{R})$,
der Hilbertraum ist $L^2(\Gamma\backslash\mathbb{H})$, und der Laplace-Beltrami-Operator
wird eindeutig charakterisiert als der Differentialoperator zweiter Ordnung, der invariant
unter $\Gamma$ ist (bis auf Skalare) (\S\ref{sec:selberg}). Im Riemannschen Fall stammt die Primverschiebungsfamilie
$\{D(r) : r \in (0,2)\}$ von keiner natürlichen endlich-dimensionalen Lie-Gruppen-Operation
der durch (N1) geforderten Form ab; Theorem~\ref{thm:NE-B} etabliert,
dass kein nicht-skalarer symmetrischer Operator überhaupt mit allen $D_N(r)$ an
den getesteten Trunkierungen kommutiert. Der Prime-Hub-Graph, falls er konstruiert
würde, müsste eine explizite $G_{\mathrm{prime}}$ liefern, was Teil von
Offenem Problem~5 ist.

\subsection{Warum kein einheitliches Theorem über alle Zeta-Typ-Familien?}

Eine plausible naive Hoffnung ist, dass $\HBL$ \emph{einheitlich}
$\mathrm{YES}$ über Zeta-Typ-Familien hinweg wäre: Jede $L$-Funktion würde ihren eigenen
natürlichen $H_X$ tragen. Das Randtheorem (boundary theorem) unten ist eine präzise
Quantifizierung des Fehlschlagens dieser Hoffnung: Selbst die drei kanonischen
Instanzen $\{\zeta, Z_\Gamma, G_\mathrm{prime}\}$ realisieren bereits alle drei
Werte $\{\mathrm{NO}, \mathrm{YES}, \mathrm{OPEN}\}$. Jeder einheitliche
Anspruch über alle Zeta-Typ-Familien ist daher falsch.

Eine subtilere Hoffnung ist, dass $\HBL = \mathrm{YES}$ \emph{generisch} wäre,
wobei der Riemannsche Fall auf irgendeine Weise außergewöhnlich ist. Theorem~\ref{thm:NE-B}
verbietet sogar diese Lesart an den getesteten Galerkin-Trunkierungen: Das
Fehlen einer gemeinsamen Symmetrie unter Primverschiebungsoperatoren ist ein explizites
strukturelles Merkmal der multiplikativen Primzahlstruktur, keine Pathologie. Umgekehrt
ist das Vorhandensein eines Lie-Gruppen-invarianten Laplace-Operators in Selbergs
Setting eine Folge der spezifischen Geometrie von $\Gamma\backslash\mathbb{H}$ und
nicht per se der Zeta-Analytizität. Die Branch-Lokalitäts-These lautet, dass
der $\HBL$-Status vom \emph{algebraischen Substrat} der Familie determiniert wird,
unabhängig von deren zeta-analytischen Merkmalen.

\begin{remark}[Relativer Status]\label{rem:relative-hbl}
Die Notation $\HBL(X)$ ist immer relativ zur deklarierten
Transfer-Präsentation $\mathcal{T}_X$ aus Definition~\ref{def:zeta-family}
zu lesen. Wo Mehrdeutigkeit droht, bedeutet sie
$\HBL_{\mathcal{T}}(X)$. Die Riemann-Zeile unten ist also eine Aussage
über die klassische Primverschiebungs-Transferpräsentation. Connes'
adelisches Programm, das unten diskutiert wird, ist eine andere
Transferpräsentation und bleibt eine mögliche offene Route zu einem
adelischen $\HBL_{\mathcal{T}}(\zeta)=\mathrm{YES}$.
\end{remark}

\begin{remark}[Striktes und resonantes YES]\label{rem:yes-types}
Wir unterscheiden zwei positive Statuswerte. $\mathrm{YES}_{\mathrm{sa}}$
bezeichnet die strikte selbstadjungierte Punkt-Spektrum-Realisierung aus
Definition~\ref{def:hp-bl}. $\mathrm{YES}_{\mathrm{res}}$ bezeichnet das
nicht-kompakte Flow-Zeta-Analogon, in dem ein natürlicher selbstadjungierter
Operator ein Streu-/Resolventenproblem kontrolliert und die Nullstellen
Resonanzen statt gewöhnlicher Eigenwerte sind. Das ist für die Zoo-Taxonomie
nützlich, aber ohne diese Resonanzqualifikation keine wörtliche Instanz von
(HP1).
\end{remark}

\subsection{Eine Drei-Klassen-Taxonomie des Zeta-Zoos}
\label{sec:three-class-taxonomy}

Die Branch-Lokalitäts-Trichotomie $\HBL \in \{\mathrm{YES}, \mathrm{NO}, \mathrm{OPEN}\}$
aus Definition~\ref{def:hpbl-status} klassifiziert Familien nach der
\emph{Existenz} eines natürlichen Hilbert-P\'olya-Operators. Eine zweite,
weitgehend orthogonale Klassifizierung erfasst die \emph{Art des Spektrums}, das
ein solcher Operator, falls er existiert, zulässt. Drei Klassen treten hervor:

\begin{center}
\small
\setlength{\tabcolsep}{3pt}
\renewcommand{\arraystretch}{1.18}
\begin{tabular}{@{}p{0.19\linewidth} p{0.18\linewidth} p{0.12\linewidth} p{0.18\linewidth} p{0.17\linewidth}@{}}
\toprule
\textbf{Klasse} & \textbf{Transfer-Algebra} & \textbf{Typischer HP} & \textbf{Spektrum} & \textbf{Zoo-Beispiele} \\
\midrule
Arithmetische Zetas & Prim/Primideal-Halbgruppe & NO/SGE-NO & Diskret (Nullstellen) & Riemann, Dirichlet, Dedekind \\
Geometr. Zetas (kompakt) & Diskretes $\Gamma$ auf kompaktem Quotient & YES$_{\mathrm{sa}}$ & Diskret (Laplace) & Selberg, Ihara \\
Flow-Zetas (nicht-kompakt) & Kont. Lie-Gruppe auf nicht-kompaktem Raum & YES$_{\mathrm{res}}{}^*$ & Kontinuierlich + Resonanzen & CRM, Ruelle-Anosov, BH-QNM \\
\bottomrule
\end{tabular}
\end{center}
${}^*$ realisiert im Resonanzsinne: Der HP-Operator (Casimir der
operierenden Lie-Gruppe) besitzt absolut stetiges Spektrum, und die Zeta-Nullstellen
entsprechen Streuresonanzen (Scattering Resonances) in der unteren Halbebene (vgl. \cite{DyatlovZworski2019}).

Die drei Klassen sind orthogonal zur $\HBL$-Trichotomie: Selberg und CRM
sitzen auf positiven Seiten der SGE-Dichotomie, aber Selberg ist
$\mathrm{YES}_{\mathrm{sa}}$, während CRM nur
$\mathrm{YES}_{\mathrm{res}}$ ist. Riemann und Dedekind teilen sich die arithmetische Klasse,
bleiben aber von den geometrischen und Flow-Klassen vollständig getrennt. Wir werden eine kanonische
Vertretung für jede Klasse in den Abschnitten \ref{sec:riemann}--\ref{sec:crm-flow} aufzeigen
und geben das vollständige Randtheorem in Abschnitt~\ref{sec:boundary}.

\subsection{Beziehung zu Connes' adelischem Framework}

In Connes' Programm \cite{Connes1999, Connes2026} ist der relevante
"Raum", auf dem ein vermeintlicher $H_\zeta$ operiert, ein adelischer
Quotient $\mathbb{A}_K / K^\times$, und die natürliche Symmetrie ist die
$\mathrm{GL}_1$-Operation der Ideleklassen. Unsere Formulierung in Definition~\ref{def:natural-hp}
schließt diese Architektur nicht aus: Wenn Connes' Programm erfolgreich ist,
würde sein Operator (N1)--(N3) mit $G_X = \mathrm{GL}_1(\mathbb{A}_K/K^\times)$
erfüllen und ein explizites $\HBL = \mathrm{YES}$-Zertifikat für $\zeta(s)$
im adelischen Hilbertraum darstellen. Diese Möglichkeit ist derzeit offen:
Sie würde auf endlichen Dimensionen mit Theorem~\ref{thm:NE-B} koexistieren,
weil das adelische natürliche $G_X$ in der Galerkin-Trunkierung nicht sichtbar ist.
Unser Randtheorem betrifft daher die Hilbert-P\'olya-Struktur im Hinblick auf die
\emph{Primverschiebungs-}Transfer-Operatoren $\mathcal{T}_\zeta$, was die relevante
Struktur für die Weilsche quadratische Form und das Even-Dominance-Beweisprogramm
aus \cite{Geiger2026RHII} ist.

Die folgenden Abschnitte stellen jeweils eine kanonische Instanz jedes
$\HBL$-Modus vor.

% ============================================================
\section{Instanz A: Riemann (NE-A und NE-B)}
\label{sec:riemann}

\subsection{Nichtexistenz absolut totaler Positivität (NE-A)}
\label{sec:ne-a}

Wir zitieren das folgende Theorem aus \cite{Geiger2026RHII}:

\begin{theorem}[NE-A, {\cite[Thm.~NE-A]{Geiger2026RHII}}]\label{thm:NE-A}
Der Fouriermultiplikator
$M(\xi) = -2\,\Re[\zeta'/\zeta(\halfline + i\xi)]$
des Primverschiebungsoperators nimmt auf einer Menge positiven Lebesgue-Maßes
in jedem Intervall $[-L, L]$ mit $L > \gamma_1$ (der ersten positiven Zeta-Ordinate)
negative Werte an. Folglich ist der Primverschiebungsoperator
nicht absolut total-positiv ($\mathrm{PF}_\infty$).
\end{theorem}

\begin{remark}\label{rem:ne-a-empirical}
Empirisch (vgl. \cite[\S4.2]{Geiger2026RHII}) liefert die Abtastung von $M(\xi)$
auf einem gleichmäßigen Gitter von $10^4$ Punkten in $\xi \in [0,100]$ unter Ausschluss
von $\varepsilon$-Umgebungen der ersten 29 Ordinaten ca. $49\%$ Punkte mit $M(\xi) < 0$.
Dies illustriert das Theorem; es ersetzt nicht das strukturelle Argument.
\end{remark}

\subsection{Kein universeller kommutierender Operator (NE-B)}
\label{sec:ne-b}

\begin{theorem}[NE-B endliches Zertifikat, {\cite[Thm.~NE-B]{Geiger2026RHII}}]\label{thm:NE-B}
Für jede registrierte Galerkin-Trunkierung $N \leq 15$ und die in
\cite{Geiger2026RHII} berichteten zertifizierten dichten Stichproben
normalisierter Verschiebungen $r_1, \ldots, r_k \in (0,2)$ ist der
berechnete gemeinsame symmetrische Zentralisator der
Primverschiebungsmatrizen $D_N(r_j)$ eindimensional. Äquivalent ist
innerhalb dieser endlichen zertifizierten Trunkierungen jede simultan
kommutierende symmetrische Matrix ein skalares Vielfaches der
Einheitsmatrix.
\end{theorem}

\begin{conjecture}[NE-B-full, {\cite[Conj.~\ref{thm:NE-B}-full]{Geiger2026RHII}}]\label{conj:NE-B-full}
Theorem~\ref{thm:NE-B} dehnt sich auf $N = \infty$ aus: Es existiert kein nicht-trivialer
universeller kommutierender Operator für die Riemannsche Primverschiebungs-Familie.
\end{conjecture}

\subsection{Konsequenz für \texorpdfstring{$\HBL(\zeta)$}{HP-BL(zeta)}}

Theoreme~\ref{thm:NE-A} und~\ref{thm:NE-B} schließen zusammen die beiden
natürlichsten Routen zu (HP1)--(HP3) für die Riemann-Zeta-Familie aus.
Kein absolut total-positiver Primverschiebungs-Operator kann als $H_\zeta$ dienen
(Route A), und es existiert kein universeller kommutierender selbstadjungierter
symmetrischer Operator bei irgendeiner getesteten Galerkin-Trunkierung (Route B). Unter
Vermutung~\ref{conj:NE-B-full} gilt $\HBL(\zeta) = \mathrm{NO}$.

\begin{remark}[Paper IIs bedingte RH-Route erfordert keine HP-Struktur]
\label{rem:riemann-without-hp}
Die Even-Dominance-Route zur Riemannschen Vermutung in
\cite{Geiger2026RHII} wird derzeit als bedingte Reduktions- und
Diagnoseroute behandelt: Sie verläuft über primzahlgewichtete Summation
(Shift-Parity-Lemma) plus Frontier-Dominanz-Argumente unter Verwendung
der Weilschen quadratischen Form, aber die endliche CAP-Evidenz ist für
sich genommen kein unbedingter RH-Beweis. Keiner dieser Schritte
benötigt (HP1)--(HP3). Dies ist konsistent mit der
Branch-Lokalitäts-These: Selbst wenn $\HBL(\zeta) = \mathrm{NO}$
(bedingt durch Vermutung~\ref{conj:NE-B-full}), kann die
Weil-Form-Architektur die relevante Nicht-HP-Route bleiben, vorbehaltlich
ihres eigenen Odd-Sector-/AVG-Abschlusses.
\end{remark}

\begin{remark}[Dokumentationsstatus der Direct-Frontier-Route]
\label{rem:riemann-critique-review}
Wir zitieren \cite{Geiger2026RHII} als Companion-Programm, das die
Even-Dominance- / Direct-Frontier-Reduktionsroute dokumentiert. Eine
systematische Kritik der Trilogie
\cite{Geiger2026RHII,Geiger2026RHIII} --- welche (K1) die globale
Extrapolation von der $\lambda = 100$ CAP-Baseline, (K2) die Transparenz
der PNT/Mertens-Konstanten $C_{\mathrm{shift}}, C_{\mathrm{norm}}$, (K3)
die Galerkin-Konvergenz zum unendlich-dimensionalen Weil-Operator, (K4)
die Implikation-versus-Äquivalenz-Struktur der Reduktionskette A1--A8,
(K5) die präzise Charakterisierung des Form-Definitionsbereichs von
$Q_\zeta$ und (K6) die gesamte Prüfer-Ergonomie abdeckte --- wird hier
als Liste der verbleibenden Guardrails dieser Route behandelt. Im
aktuellen Status sind diese Guardrails nicht kosmetisch: Die öffentliche
RH-Folgerung bleibt bedingt, bis der Odd-Sector-Lower-Bound /
AVG-artige Abschluss und die relevante Galerkin-zu-Weil-Brücke vorliegen.
Diese Qualifikation berührt weder das Randtheorem noch die
SGE-Klassifizierung des vorliegenden Papers; sie verhindert nur, dass das
Companion-Programm hier als unbedingter RH-Beweis verwendet wird.
\end{remark}

% ============================================================
\section{Instanz B: Selberg (Casimir-Operator)}
\label{sec:selberg}

\subsection{Der Laplace-Beltrami als expliziter Hilbert-P\'olya-Operator}

Sei $\Gamma \subset \PSL_2(\mathbb{R})$ eine ko-kompakte Fuchssche Gruppe
und $\mathcal{F} = \Gamma\backslash\mathbb{H}$ die zugehörige kompakte
hyperbolische Fläche. Die Selbergsche Zetafunktion $Z_\Gamma(s)$ hat
nicht-triviale Nullstellen $\rho_n = \halfline + ir_n$, die in Beziehung stehen
zum Spektrum des Laplace-Beltrami-Operators $\Delta_\mathcal{F}$ auf
$L^2(\mathcal{F})$ via
\begin{equation}
\label{eq:selberg-spectral}
  \Delta_\mathcal{F}\,\phi_n = \lambda_n\phi_n, \qquad
  \lambda_n = \tfrac{1}{4} + r_n^2,
\end{equation}
mit endlich vielen Ausnahmen durch kleine Eigenwerte \cite{Selberg1956}.

Für diese Zeile bezeichnet $\mathcal{T}_\Gamma$ die Spurformel-/
Wave-Kernel-Algebra, die durch den Funktionalkalkül von
$\sqrt{\Delta_\mathcal{F}-1/4}$ erzeugt wird; geschlossene Geodäten gehen
über die Selbergsche Spurformel ein. Wir behaupten keine Kommutation mit
beliebigen Ruelle-, symbolischen oder gewichteten Transferoperatoren auf
anisotropen Räumen.

\begin{theorem}[Selberg]\label{thm:selberg-hp}
Nach Abtrennung des standardmäßigen endlichen Ausnahmesektors
$\{\lambda_n<1/4\}$ erfüllt der Operator
$H_{Z_\Gamma}:=\sqrt{\Delta_\mathcal{F}-1/4}$ auf dem Spektralunterraum
$\lambda_n\geq 1/4$ die Bedingungen (HP1)--(HP3) aus
Definition~\ref{def:hp-bl} für den prinzipalen Kritische-Gerade-Anteil
der Zeta-Typ-Familie
$(Z_\Gamma, \mathcal{T}_\Gamma, \mathcal{E}_\Gamma)$. Die
außergewöhnlichen reellen Nullstellen sind eine endliche Selberg-Korrektur
und nicht Teil des Hilbert-P\'olya-Ordinatenspektrums. In diesem
standardmäßig qualifizierten Sinn gilt $\HBL(Z_\Gamma)=\mathrm{YES}$.
\end{theorem}

\begin{proof}[Skizze]
(HP1) folgt aus~\eqref{eq:selberg-spectral} auf dem Unterraum
$\lambda_n\geq 1/4$, wo $r_n$ reell ist. (HP2) folgt aus
der $\Gamma$-Invarianz für die kanonische Spurformel-/Wave-Kernel-Algebra,
die durch den Isometriefluss erzeugt wird: Der Laplace-Beltrami-Operator
kommutiert mit den relevanten Isometrie-Pullbacks und ihren
Funktionalkalkül-Operatoren.
(HP3) folgt aus der Eindeutigkeit des Laplace-Operators als dem
invarianten Differentialoperator zweiter Ordnung unter $\Gamma$ (bis auf Skalare).
Da $\Gamma\backslash\mathbb{H}$ in dieser Instanz kompakt ist, hat
$\Delta_\mathcal{F}$ diskretes Spektrum; die endlich vielen kleinen
Eigenwerte werden als die üblichen außergewöhnlichen Selberg-Nullstellen
behandelt, nicht über einen absolut stetigen Spektralanteil.
\end{proof}

\subsection{Selbergs Spurformel: der explizite Mechanismus}
\label{sec:selberg-trace}

Die Selbergsche Spurformel \cite{Selberg1956} setzt schematisch spektrale und geometrische
Daten in Beziehung:
\begin{equation}
\label{eq:selberg-trace}
  \sum_n \widehat{h}(r_n)
  = \frac{\operatorname{vol}(\mathcal{F})}{4\pi} \int h(r)\,r\tanh(\pi r)\,dr
  + \sum_{\{\gamma\}} \frac{\ell(\gamma_0)}{2\sinh(\ell(\gamma)/2)}\,
      g(\ell(\gamma)),
\end{equation}
wobei $r_n$ die spektralen Parameter von $\Delta_\mathcal{F}$ sind, die
zweite Summe über primitive geschlossene Geodäten $\{\gamma_0\}$ und deren
Iterierte läuft, und $h, g$ ein durch Fouriertransformation zusammenhängendes
Standard-Testfunktionspaar sind. Die geodätische Seite von \eqref{eq:selberg-trace}
ist das \emph{multiplikativ-geometrische Analogon} zur Prim-Seite der expliziten Formel
von Riemann--Weil.

Die $\Gamma$-Invarianz von $\Delta_\mathcal{F}$ ist der
\emph{strukturelle Grund}, warum die geodätische Seite diese geschlossene Form hat:
Geodäten $\gamma$ sind Orbits unter der $\PSL_2(\mathbb{R})$-Operation
modulo $\Gamma$, und die Transferoperatoren $T_{\ell_\gamma}$ sind
konjugiert zu $\Gamma$-Translationen entlang $\gamma$. Der Casimir-Operator
der einbettenden $\PSL_2(\mathbb{R})$ wirkt als Erzeuger zweiter Ordnung dieser geodätischen
Flüsse und kommutiert daher simultan mit jedem $T_{\ell_\gamma}$. Dies ist
die konkrete Verifizierung von (N1)--(N3) aus Definition~\ref{def:natural-hp}.

\subsection{Was Selberg von Riemann unterscheidet}
\label{sec:selberg-vs-riemann}

Die Riemannsche Zeta-Familie lässt ein formales Analogon zu \eqref{eq:selberg-trace}
in Form der Weilschen expliziten Formel zu:
\begin{equation}
\label{eq:weil-riemann}
  \sum_k h(\gamma_k^\zeta)
  = W_\infty[h]
  - 2\sum_{p,m} \frac{\log p}{p^{m/2}}\,\widehat{h}(m\log p),
\end{equation}
mit archimedischem Gewicht $W_\infty$, das den Gamma-Faktor kodiert. Die
Prim-Seite von \eqref{eq:weil-riemann} ist \emph{multiplikativ}: Sie ist
durch die Halbgruppe der Primzahlpotenzen indiziert, nicht durch Orbits einer
Lie-Gruppen-Operation. Es gibt keine natürliche kontinuierliche Gruppe,
deren diskrete Orbits $\{p^m\}$ reproduzieren, und Theorem~\ref{thm:NE-B} bestätigt
dies strukturell: Kein nicht-skalarer symmetrischer Operator kommutiert mit
den Primverschiebungs-Operatoren bei den getesteten Galerkin-Trunkierungen.

Der genaue strukturelle Unterschied zwischen \eqref{eq:selberg-trace} und
\eqref{eq:weil-riemann} lautet daher:
\begin{center}
\begin{tabular}{@{}lll@{}}
\toprule
Merkmal & Selberg (HP-BL YES) & Riemann (HP-BL NO${}^*$) \\
\midrule
Index-Menge & geschlossene Geodäten $\{\gamma\}$ &
  Primzahlpotenzen $\{p^m\}$ \\
Gruppen-Operation & $\PSL_2(\mathbb{R})$ auf $\mathbb{H}$ & keine (NE-B) \\
Kommut. Hamilton-Op. & $\Delta_\mathcal{F}$ = Casimir & keiner für $N \leq 15$ \\
Längen-Struktur & kontinuierlich $\ell(\gamma) \in \mathbb{R}_{>0}$ &
  diskret $\log p \in \mathbb{R}_{>0}$ \\
\bottomrule
\end{tabular}
\end{center}
Die multiplikative Arithmetik der Primzahlen ist \emph{nicht die Orbit-Struktur
irgendeiner natürlichen kontinuierlichen Symmetriegruppe}. Dies ist der
Inhalt der Branch-Lokalität: Das algebraische Substrat der
Transfer-Operator-Familie entscheidet über den HP-BL-Status, unabhängig
von der zeta-analytischen Ähnlichkeit zwischen \eqref{eq:selberg-trace}
und \eqref{eq:weil-riemann}.

% ============================================================
\section{Instanz C: Prime-Hub (Offenes Problem 5)}
\label{sec:primehub}

Offenes Problem 5 aus \cite{Geiger2026RHIII}, im Folgenden OP5, fordert einen
expliziten topologisch mischenden gerichteten Graphen
$G_\mathrm{prime} = (V, E)$ (oder äquivalent einen Fluss oder ein dynamisches System),
der fünf Bedingungen erfüllt:
\begin{enumerate}
\item[(OP5.1)] Primitive periodische Orbits $\{\gamma_p\}$ stehen in
  Bijektion zu den Primzahlen.
\item[(OP5.2)] $\ell(\gamma_p) = \log p$ für jede Primzahl $p$.
\item[(OP5.3)] Der zugehörige Transfer-Operator $\mathcal{L}_s$ hat
  die Fredholm-Determinante $\det(I - \mathcal{L}_s)^{\mathrm{vac}} =
  1/\zeta(s)$.
\item[(OP5.4)] Der Spektralradius von $\mathcal{L}_s$ ist kleiner als $1$
  für $\Re(s) > 1/2$.
\item[(OP5.5)] Der Eigenwert $1$ tritt genau an den nicht-trivialen Nullstellen
  von $\zeta(s)$ auf.
\end{enumerate}
Keine Konstruktion, die (OP5.1)--(OP5.5) simultan erfüllt, ist
seit der informellen Hilbert-P\'olya-Anregung von 1914 bekannt
\cite{Polya1982letter}. Das Programm wird in \cite{Connes2026}
besprochen und übersichtlich dargestellt; Connes' \emph{Letter Through Time}
(ebd., letzte Abschnitte) schlägt eine auf der Spurformel basierende
Konvergenzstrategie von endlichen zu unendlichen Euler-Produkten vor, jedoch wird ein konstruktiver
$G_\mathrm{prime}$ nicht vorgezeigt.

\subsection{Hindernis C1: Weyl-Gesetz-Inkompatibilität}
\label{sec:obs-C1}

\begin{proposition}[Weyl-Gesetz-Inkompatibilität]\label{prop:weyl-obstruction}
Sei $H$ der Laplace-Beltrami-Operator einer glatten kompakten
$d$-dimensionalen Riemannschen Mannigfaltigkeit. Dann erfüllt seine Eigenwertzählfunktion
das Weylsche Gesetz
\begin{equation}
\label{eq:weyl-classical}
  N_H(T) := \#\{\lambda \in \Spec(H) : \lambda \leq T^2\}
  \;\sim\; \frac{\omega_d \operatorname{vol}(M)}{(2\pi)^d}\,T^d
  \qquad (T \to \infty),
\end{equation}
wobei $\omega_d$ das Volumen der Einheitskugel in $\mathbb{R}^d$ ist. Die
Zählfunktion der Riemann-Nullstellen lautet
\begin{equation}
\label{eq:riemann-N}
  N_\zeta(T) := \#\{\gamma_k^\zeta : 0 < \gamma_k^\zeta \leq T\}
  =
  \frac{T}{2\pi}\log\!\frac{T}{2\pi e} + O(\log T)
\end{equation}
(klassisches Resultat von Riemann-von Mangoldt, vgl.\ \cite[\S9.4]{Titchmarsh1986}).
Die asymptotischen Raten \eqref{eq:weyl-classical} und
\eqref{eq:riemann-N} sind für jede ganze Zahl $d \geq 1$ inkompatibel.
Insbesondere kann kein klassischer geometrischer Laplace-Operator auf einer glatten kompakten
Mannigfaltigkeit $\Spec(H_\zeta)$ realisieren.
\end{proposition}

\begin{proof}
Das Potenzgesetz \eqref{eq:weyl-classical} steht im Kontrast zur
logarithmischen Dichte \eqref{eq:riemann-N}: für jedes $d \geq 1$ gilt für
das Verhältnis $N_\zeta(T) / N_H(T) \to 0$ (falls $d \geq 2$) oder $\to \infty$ für
$T \to \infty$ (falls $d = 1$), also macht keine Wahl von $d$ und $\operatorname{vol}(M)$
\eqref{eq:weyl-classical} mit \eqref{eq:riemann-N} passend.
Folglich kann kein Kandidat $H_\zeta$ mit
$\Spec(H_\zeta) = \{\gamma_k^\zeta\}$ der Laplace-Operator einer
glatten kompakten Riemannschen Mannigfaltigkeit sein.
\end{proof}

\begin{remark}\label{rem:weyl-workarounds}
Proposition~\ref{prop:weyl-obstruction} schließt glatte kompakte
Mannigfaltigkeits-Laplace-Operatoren aus, jedoch nicht alle Möglichkeiten: nicht-kompakte Spektren,
hyperbolisch-ähnliches Volumenwachstum (Berrys chaotische Quantisierung
\cite{Berry1985, BerryKeating1999}) oder adelische Räume
\cite{Connes1999, Connes2026} bleiben als Kandidatenarchitekturen bestehen. Das
Hindernis ist daher eine Beschränkung der \emph{Art} des
geometrischen Substrats.
\end{remark}

\subsection{Hindernis C2: Spurklasse-Regularität}
\label{sec:obs-C2}

\begin{proposition}[Spurklasse-Regularität]\label{prop:trace-class}
Die Fredholm-Determinante $\det(I - \mathcal{L}_s)^{\mathrm{vac}}$ in
Bedingung (OP5.3) ist nur dann wohldefiniert, wenn $\mathcal{L}_s$ eine
Spurklasse auf seinem umgebenden Hilbertraum ist. Für einen Graphentransferoperator,
dessen primitive Orbits logarithmische Längen
$\ell(\gamma_p) = \log p$ haben, muss die Summierbarkeitsbedingung
\begin{equation}
\label{eq:trace-class-condition}
  \sum_{p,m} \frac{1}{p^{m/2}} \cdot |\lambda_p^m|
  \;<\; \infty
\end{equation}
(wobei $\lambda_p$ der Eigenwert von $\mathcal{L}_s$ entlang des primitiven
Zyklus $\gamma_p$ ist) gleichmäßig in dem Bereich gelten, in dem (OP5.4)
einen Spektralradius kleiner als $1$ verlangt. Keine explizite Konstruktion,
die (OP5.1)--(OP5.5) und
\eqref{eq:trace-class-condition} simultan erfüllt, wurde bisher vorgezeigt; dies ist
im Geiger-Baum als Hindernis K4 in \cite{Geiger2026RHI} bekannt.
\end{proposition}

\begin{remark}\label{rem:trace-class-connection}
Die Schwierigkeit hinter Proposition~\ref{prop:trace-class} ist, dass
(OP5.4) und \eqref{eq:trace-class-condition} in entgegengesetzte
Richtungen ziehen: (OP5.4) beschränkt die Eigenwerte $\lambda_p$ von oben
durch $1$, während \eqref{eq:trace-class-condition} ihre
\emph{Summe} als absolut konvergent erzwingt. Für einen „generischen“ Graphen
mit unendlich vielen primitiven Orbits liefern natürliche Operatormodelle entweder
eine Verletzung von (OP5.4) (Eigenwerte häufen sich nahe $1$) oder
eine Verletzung von \eqref{eq:trace-class-condition} (langsamer Abfall). Dieses
Nadelöhr zu passieren, ist der substanzielle Inhalt von OP5.
\end{remark}

\subsection{Hindernis C3: NE-A-Kompatibilität}
\label{sec:obs-C3}

\begin{proposition}[NE-A-Kompatibilität]\label{prop:ne-a-compat}
Jeder Hilbert-P\'olya-Kandidat $H_\mathrm{prime}$ für $\zeta$ determiniert implizit
den Fouriermultiplikator
$M(\xi) = -2\,\Re[\zeta'/\zeta(1/2 + i\xi)]$ des Primverschiebungsoperators
via die Spurformel von $(H_\mathrm{prime}, \mathcal{L}_s)$.
Nach Theorem~\ref{thm:NE-A} ist $M(\xi) < 0$ auf einer Menge von positivem
Lebesgue-Maß. Folglich kann $H_\mathrm{prime}$ nicht als ein
nicht-negativer selbstadjungierter Operator mit nicht-negativem Fourier-Symbol
auf glatten Geometrien entstehen; er muss intrinsische Vorzeichenstruktur
in seiner Spektraltransformation tragen.
\end{proposition}

\begin{proof}[Skizze]
Ein Hilbert-P\'olya-Kandidat muss die Weilsche explizite Formel
(Abschnitt~\ref{sec:weil-qf}) via seiner Spur reproduzieren. Der Kern der
Spurformel ist genau $M(\xi)$. Standard-Laplace-Operatoren auf glatter Geometrie
erzeugen nicht-negative Fouriermultiplikatoren; das Vorzeichen von $M(\xi)$ auf einer
Menge mit positivem Maß schließt folglich derartige Laplace-Operatoren aus.
Eine Konstruktion muss Strukturen mit inhärentem Vorzeichenwechsel (z. B.
pseudodifferentielle oder nichtkommutative adelische) verwenden.
\end{proof}

\begin{remark}\label{rem:ne-a-connection}
Die Hindernisse C1 (glatte Geometrie ausgeschlossen) und C3 (nicht-negative
Symbole ausgeschlossen) sind unabhängig, verstärken sich aber gegenseitig: C1
schließt Mannigfaltigkeitsarchitekturen via Dichte aus, C3 schließt sie
via Symbol-Vorzeichen aus. Selbst wenn ein Workaround (z.B.
hyperbolische chaotische Quantisierung) C1 umgeht, muss er immer noch
ein Symbol erzeugen, das C3 erfüllt.
\end{remark}

\subsection{Hindernis C4: Determinantenorientierung}
\label{sec:obs-C4}

\begin{proposition}[Fredholm-Orientierung]\label{prop:fredholm-orientation}
Sei $U$ eine Umgebung einer nichttrivialen Nullstelle $\rho$ von
$\zeta$, und sei $s\mapsto\mathcal{L}_s$ eine holomorphe
Spurklasse-Familie auf $U$. Wenn
\[
  D(s):=\det(I-\mathcal{L}_s)
\]
eine gewöhnliche Fredholm-Determinante ist und (OP5.3) nach
meromorpher Fortsetzung als $D(s)=1/\zeta(s)$ nahe $\rho$ gelesen
wird, dann kann $\rho$ kein gewöhnliches Eigenwert-eins-Ereignis von
$\mathcal{L}_s$ sein. Denn ein Eigenwert-eins-Ereignis erzwingt
$D(\rho)=0$, während $1/\zeta(s)$ bei $\rho$ einen Pol hat.
\end{proposition}

\begin{proof}
Für eine holomorphe Spurklasse-Familie ist die Fredholm-Determinante
holomorph. Wenn $1$ in das Spektrum von $\mathcal{L}_\rho$ eintritt,
ist $I-\mathcal{L}_\rho$ nicht invertierbar und die gewöhnliche
Fredholm-Determinante verschwindet bei $\rho$. Ist $\rho$ dagegen eine
Nullstelle von $\zeta$ der Ordnung $m\geq 1$, dann hat $1/\zeta(s)$
bei $\rho$ einen Pol der Ordnung $m$. Eine holomorphe
Determinantennullstelle und ein meromorpher Pol können nicht dasselbe
lokale Ereignis sein. Daher können die beiden OP5-Signale nicht in
einer einzigen ungradierten gewöhnlichen Fredholm-Determinante liegen.
\end{proof}

\begin{remark}[Gradierte Prime-Hub-Lesart]\label{rem:graded-primehub}
Das Hindernis widerlegt OP5 nicht; es präzisiert seine Semantik. Das
Symbol ``vac'' in (OP5.3) muss eine Quotienten-, regularisierte oder
gradierte Determinante codieren, oder OP5 muss in zwei Kanäle zerlegt
werden:
\[
  D_E(s)=\prod_p(1-p^{-s})=\frac{1}{\zeta(s)},\qquad
  D_Z(s)=\frac{\xi(s)}{\xi(0)},\qquad
  G(s)=\frac{\xi(s)}{\xi(0)\zeta(s)}
\]
mit
\[
  D_Z(s)=\frac{G(s)}{D_E(s)}.
\]
Der Euler-Kanal $D_E$ trägt Prim-Orbits und Möbius-Parität; der
Spektralkanal $D_Z$ trägt das Signal der nichttrivialen Nullstellen;
und $G$ enthält die Vervollständigungsfaktoren. Äquivalent kann man
einen $\mathbb{Z}_2$-gradierten Transferkomplex suchen mit
\[
  \operatorname{Sdet}(I-\mathcal{L}_s)
  =\frac{\det(I-\mathcal{L}_s^+)}
          {\det(I-\mathcal{L}_s^-)}
  =\frac{1}{\zeta(s)},
\]
sodass Eigenwert-eins-Ereignisse an Zeta-Nullstellen im Nennersektor
auftreten. Das elementare endliche Modell ist der Möbius-gradierte
Primoperator $A_s e_p=p^{-s}e_p$, für den
\[
  \det(I-A_s)=\operatorname{Str}_{\Lambda^\ast\ell^2(P)}(\Lambda A_s)
  =\sum_{n\ \mathrm{squarefree}}\mu(n)n^{-s}
  =\prod_p(1-p^{-s})
\]
für $\Re(s)>1$ in der unendlichen Primmenge gilt. Dies ist in der
Arbeitsnotiz \cite{MobiusGradedPrimeHub} festgehalten. Die aktuelle
numerische Fortsetzung \cite{PrimeWeilBridgeDefect} zeigt, warum ein
nicht-tautologischer Prime--Weil-Defekt jede direkte Einfügung einer
Nullstellenliste ersetzen muss.
\end{remark}

\subsection{Konsequenz}

Die Kombination der Propositionen~\ref{prop:weyl-obstruction},
\ref{prop:trace-class}, \ref{prop:ne-a-compat} und
\ref{prop:fredholm-orientation} identifiziert vier
strukturelle Hindernisse, durch die jede explizite Prime-Hub-Konstruktion
navigieren muss. Keines der vier ist für sich genommen ein Unmöglichkeitsbeweis;
zusammen grenzen sie den engen Parameterraum ein, in dem
ein gültiger $G_\mathrm{prime}$ existieren könnte. Folglich
gilt $\HBL(G_\mathrm{prime}) = \mathrm{OPEN}$, mit den Hindernissen C1--C4
als explizite Wegweiser für jede zukünftige Konstruktion. Connes'
\emph{Letter Through Time} \cite{Connes2026} exemplifiziert die Art von
nicht-glatter, nichtkommutativer Architektur, die prinzipiell
alle vier durchsteuern könnte; ob das Programm abschließt, bleibt offen.

% ============================================================
\section{Instanz D: CRM als Flow-Zeta (FST-Kosmologie Brücke)}
\label{sec:crm-flow}

Die Instanzen A, B, C repräsentieren die arithmetischen und geometrisch-kompakten
Klassen der Drei-Klassen-Taxonomie aus Abschnitt~\ref{sec:three-class-taxonomy}.
Die Flow-Zeta-Klasse ist durch ein konkretes Objekt der
physikalischen Seite der FST bevölkert: dem Kooperativen Renormierungsmodell (CRM)
der kosmologischen Dunklen Energie \cite{RFEPFramework}. Seine mathematische
Struktur ist detailliert in \cite{CRMFlowZetaNote} entwickelt; wir
fassen hier die Kernpunkte zusammen.

\subsection{CRM als eine einparametrige Lie-Gruppe}

Die CRM-Kompositionsregel, die die Dunkle-Energie-Evolution steuert, lautet
\begin{equation}
\label{eq:crm-composition}
  x \oplus y = \frac{x+y}{1+xy}, \qquad x, y \in (-1, 1).
\end{equation}
Setzt man $x = \tanh u$, $y = \tanh v$, so ergibt sich $x \oplus y = \tanh(u+v)$,
sodass $(x, \oplus)$ isomorph zur additiven Gruppe $(\mathbb{R},
+)$ ist. Die Komposition ist daher die Lorentz-Boost-Addition, und
die relevante Transfer-Algebra ist eine \emph{einparametrige Lie-Gruppe},
die als Untergruppe von $\PSL_2(\mathbb{R})$ realisiert wird und auf der
hyperbolischen Ebene $\mathbb{H}$ operiert.

\begin{proposition}[CRM Transfer-Algebra]\label{prop:crm-transfer}
Für die CRM-Familie $X = \mathrm{CRM}$ ist die Transfer-Algebra
$(I_{\mathrm{CRM}}, \cdot) = (\mathbb{R}, +)$, eine lokalkompakte
abelsche Lie-Gruppe. Folglich gilt unter Vermutung~\ref{conj:sge}
$\HBL(\mathrm{CRM}) = \mathrm{YES}_{\mathrm{res}}$ im
nicht-kompakten Resonanzsinn von Bemerkung~\ref{rem:yes-types}.
\end{proposition}

\subsection{Der natürliche Hilbert-P\'olya-Operator und sein Spektrum}

Der Casimir-Operator von $\mathfrak{sl}_2(\mathbb{R})$, eingeschränkt auf die
Boost-erzeugende Komponente, ist bis auf Äquivalenz der
Laplace-Beltrami-Operator $\Delta_{\mathbb{H}}$ auf der hyperbolischen
Ebene \cite{DyatlovZworski2019}. Im nicht-kompakten Fall erfüllt
$\Delta_{\mathbb{H}}$ nicht die strikte Punkt-Spektrum-Form von (HP1) aus
Definition~\ref{def:hp-bl}; sein absolut stetiges Spektrum ist
$[1/4,\infty)$. Was er liefert, ist die natürliche selbstadjungierte
Provenienz und Kommutation mit der vollen $\PSL_2(\mathbb{R})$-Operation,
mithin auch mit dem CRM-Fluss, die die zugehörige Resonanztheorie steuert.

Die \emph{Nullstellen} von $\zeta_{\mathrm{CRM}}$ sind in diesem Setting
\emph{Streuresonanzen}: Pole der meromorphen Fortsetzung
der Resolvente $(\Delta_{\mathbb{H}} - s(s-1))^{-1}$ durch das
wesentliche Spektrum in die untere Halbebene. Auf asymptotisch
hyperbolischen Hintergründen (z. B. de~Sitter, Schwarzschild--de~Sitter)
sind diese Resonanzen exakt die \emph{Quasinormalmoden} der
entsprechenden Wellengleichung \cite{DyatlovZworski2019}. Dies realisiert
$\HBL = \mathrm{YES}_{\mathrm{res}}$ im Resonanzsinne von
Abschnitt~\ref{sec:three-class-taxonomy}.

\subsection{Die Ruelle Flow-Zeta von CRM}

Gemäß \cite{DyatlovZworski2019} ist das angemessene analytische Objekt
die Ruelle-Dyatlov-Zworski Zetafunktion
\begin{equation}
\label{eq:crm-flow-zeta}
  \zeta_{\mathrm{CRM}}(s)
  = \prod_{\gamma \in \mathcal{P}_{\mathrm{CRM}}}
    \frac{1}{1 - e^{-s\,\ell(\gamma)}},
\end{equation}
wobei $\mathcal{P}_{\mathrm{CRM}}$ die primitiven periodischen
Orbits des CRM-Flusses auf einem asymptotisch hyperbolischen Hintergrund indiziert,
und $\ell(\gamma)$ die Orbit-Länge in der Lorentz-Boost-Metrik ist.
Die Nullstellen von $\zeta_{\mathrm{CRM}}(s)$ sind die oben beschriebenen
Resonanzen.

\subsection{CRM als eine Brückenspezies}

CRM nimmt eine herausragende Position im Zeta-Zoo ein: Es
ist simultan ein Flow-Zeta-Eintrag im mathematischen Zoo (dieser
Abschnitt) und eine Muster-A-Instanziierung (Pattern A) im physikalischen RFEP-Ökosystem
\cite{RFEPFramework} (siehe auch die narrative Klammer in
Abschnitt~\ref{sec:crm-circle}). Zusammen mit Riemann (historischer
Ursprung), BSD, Hodge und P~vs~NP bildet CRM die fünfte
\emph{Brückenspezies} des Programms --- und die erste, die die
Flow-Zeta-Spalte auf beiden Seiten der FST-Zwillinge öffnet.

\subsection{Das CRM-Sättigungstheorem als eine Klassifikations-Aussage}
\label{sec:crm-saturation-classification}

Eine strukturelle Bemerkung ist angebracht. Das Sättigungstheorem von CRM~V
\cite{CRMV} besagt, dass jedes dynamische System auf $(-1,1)$, das
vier Axiome erfüllt --- Existenz einer skalaren Relaxationsvariable, Existenz eines
dissipativen Terms, Existenz eines stabilen Sättigungswerts und
monoton glatte Dynamik --- zwingend das $\tanh$-Sättigungsprofil realisiert.
Gelesen als Charakterisierung der Flow-Zeta-Klasse, ist dies
nicht bloß eine Aussage über ein einzelnes kosmologisches Modell; es ist eine
\emph{Klassifikationsaussage} für die Klasse der IR-Asymptotiken, die
kompatibel mit funktionaler Stabilität auf dem Boost-Intervall sind.

Dies ist auf der Fluss-Seite strukturell analog zur Hurwitz-/Kompaktheits-Charakterisierung,
die im arithmetischen Zweig verwendet wird: So wie
die Weilsche quadratische Form Stabilität über Zweige hinweg detektiert,
skizziert CRM~V die zulässige asymptotische Klasse im fluss-nichtkompakten
Zweig. Insbesondere können verschiedene Kandidaten für UV-Vervollständigungen
verschiedene Selektionen aus dieser Sättigungsklasse realisieren; die
Klassifikation hängt nicht von einer individuellen Realisierung ab.

Diese Beobachtung ist aus zwei Gründen bedeutsam. Erstens fasst sie CRM
nicht als einen maßgeschneiderten kosmologischen Vorschlag auf, sondern als den IR-Sektor eines
branch-lokalen strukturellen Theorems. Zweitens prognostiziert sie, dass jede UV-Theorie,
deren IR-Limes die vier Axiome von CRM~V erfüllt (z. B. ein
geometrisches Modell auf GUT-Ebene, eine RG-Fluss-Konstruktion oder eine
offene holografische Konstruktion), in derselben Sättigungsklasse
landet --- modulo quantitativem Parameter-Matching. Diesen
UV-/IR-Kompatibilitätstest operational zu machen, ist ein offenes Programm innerhalb
der FST-Kosmologie.

% ============================================================
\section{Das Randtheorem (Boundary Theorem)}
\label{sec:boundary}

\begin{theorem}[Branch-Lokalität von Hilbert-P\'olya: bedingter Kern]\label{thm:boundary}
Bedingt auf Vermutung~\ref{conj:NE-B-full} für die Riemann-Zeile und relativ
zur Primverschiebungs-Präsentation ist die $\HBL$-Eigenschaft nicht konstant
über Zeta-Typ-Familien hinweg. Kombiniert mit der Drei-Klassen-Taxonomie aus
Abschnitt~\ref{sec:three-class-taxonomy} zeigt der Zoo drei
theorem-relevante kanonische Zeilen sowie mehrere Kalibrierungszeilen. Die
folgende Tabelle ist ein Evidenz-Ledger: Nur die Riemann-, Selberg- und
Prime-Hub-Zeilen gehen in die theorem-seitige Klassifikation ein; CRM,
Dedekind und Ihara kalibrieren die Taxonomie und tragen andere Evidenzklassen.

\begin{center}
\renewcommand{\arraystretch}{1.15}
\begin{tabular}{@{}>{\raggedright\arraybackslash}p{0.21\textwidth}
                  >{\raggedright\arraybackslash}p{0.16\textwidth}
                  >{\raggedright\arraybackslash}p{0.24\textwidth}
                  >{\raggedright\arraybackslash}p{0.29\textwidth}@{}}
\toprule
Familie & Status & Reichweite / Klasse & Evidenzklasse \\
\midrule
Riemann $\zeta(s)$ & NO${}_{\mathrm{prime}}{}^*$ &
  arithmetische Primverschiebung &
  Thm.~\ref{thm:NE-A}, \ref{thm:NE-B}; Verm.~\ref{conj:NE-B-full} \\
Selberg $Z_\Gamma(s)$ & YES$_{\mathrm{sa}}$ &
  geometrisch-kompakt & Theorem-Level: Thm.~\ref{thm:selberg-hp} \\
Prime-Hub $G_\mathrm{prime}$ & OPEN &
  offene Konstruktion mit Obstruktionen C1--C4 &
  diagnostisch: Prop.~\ref{prop:weyl-obstruction}--\ref{prop:fredholm-orientation} \\
CRM (Dunkle Energie) & YES$_{\mathrm{res}}$ &
  fluss-nichtkompaktes Resonanzanalogon &
  Resonanz-Provenienz: Prop.~\ref{prop:crm-transfer}, \cite{DyatlovZworski2019} \\
Dedekind $\mathbb{Q}(\sqrt{-5})$ & SGE-prognostiziert NO &
  arithmetisch & endliche NE-B-Evidenz ($N\leq 10$):
  \S\ref{sec:sge-dedekind-test}, \cite{DedekindNEBTest} \\
Ihara $\zeta_{\text{Petersen}}$ & endlicher YES-Proxy &
  endlicher Graph & endliches Zertifikat/Proxy:
  \S\ref{sec:sge-ihara-test}, \cite{IharaPetersenTest} \\
\bottomrule
\end{tabular}
\end{center}
${}^*$ bedingt nach Vermutung~\ref{conj:NE-B-full} und relativ zur
klassischen Primverschiebungs-Präsentation von $\zeta$.
\end{theorem}

\begin{proof}
Die Riemann- und Selberg-Zeilen fassen die Abschnitte~\ref{sec:riemann}
und~\ref{sec:selberg} zusammen und liefern zwei verschiedene Statuswerte,
bedingt auf Vermutung~\ref{conj:NE-B-full} für die Riemann-Zeile. Die
Prime-Hub-Zeile aus Abschnitt~\ref{sec:primehub} zeigt, dass die
Konstruktion nicht automatisch in einen der beiden positiven/negativen Fälle
kollabiert, sondern eine eigene offene Obstruktionsklasse besitzt. Damit sind
YES$_{\mathrm{sa}}$, NO$_{\mathrm{prime}}$ und OPEN über die kanonischen
Instanzen hinweg distinkt; die $\HBL$-Eigenschaft ist also nicht-konstant.
Die CRM-, Dedekind- und Ihara-Zeilen folgen aus den späteren Kalibrierungen und
liefern keine zusätzlichen Beweisannahmen für die Nicht-Konstanz.
\end{proof}

\begin{corollary}[Implikation für die Riemannsche Vermutung]
\label{cor:rh-consequence}
Jeder unbedingte Beweis der Riemannschen Vermutung, der auf
(HP1)--(HP3) für die klassische Primverschiebungs-Präsentation von $\zeta$ beruht, muss zuerst
Vermutung~\ref{conj:NE-B-full} widerlegen. Alternativ liefern
Weil-Form-Methoden (Abschnitt~\ref{sec:weil-transversal}) eine separate
bedingte/diagnostische Companion-Route, die die Hilbert-P\'olya-Struktur
nicht benötigt; diese Route ist in \cite{Geiger2026RHII} dokumentiert,
wird hier aber nicht als unbedingter Beweis verwendet.
\end{corollary}

% ============================================================
\section{Universelle Konvexitätseindeutigkeit (UCU)}
\label{sec:ucu}

Dieser Abschnitt isoliert das Eindeutigkeitsprinzip, das in exakt
derselben abstrakten Form in allen Programmen der funktionalen Stabilitätslandschaft,
die in diesem Paper behandelt werden, operiert. Wir nennen es
\emph{Universelle Konvexitätseindeutigkeit} (UCU). Es ist das
mathematische Lemma, das dem sogenannten \emph{Pattern~A} (Muster A) in
der physikalischen Formulierung (RFEP; \cite{RFEPFramework}) und seinem
dynamischen Partner \emph{Dissipative Selektion} (DS1--DS3) zugrunde liegt.
UCU ist eine einzige technische Aussage; ihre Universalität liegt darin,
wie viele scheinbar unverbundene Objekte sie realisieren.

\subsection{Das UCU-Lemma}
\label{sec:ucu-lemma}

\begin{lemma}[UCU]\label{lem:ucu}
Sei $(\mathcal{H}, \langle\cdot,\cdot\rangle)$ ein reeller Hilbertraum,
sei $\mathcal{A} \subset \mathcal{H}$ eine abgeschlossene konvexe zulässige Klasse,
und sei $J : \mathcal{A} \to \mathbb{R}$ versehen mit:
\begin{enumerate}
\item[(C)] \emph{Strikter Konvexität:} für alle $u,v \in \mathcal{A}$
mit $u \neq v$ und alle $t \in (0,1)$ gilt
\[
  J(tu + (1-t)v) < t\,J(u) + (1-t)\,J(v);
\]
\item[(L)] \emph{Starker Unterhalbstetigkeit:} wenn $u_n \to u$
stark in $\mathcal{H}$ mit $u_n, u \in \mathcal{A}$, dann
$J(u) \leq \liminf_{n\to\infty} J(u_n)$;
\item[(K)] \emph{Koerzivität auf $\mathcal{A}$:} $J(u) \to +\infty$
wann immer $\|u\| \to \infty$ innerhalb von $\mathcal{A}$.
\end{enumerate}
Dann hat $J$ einen eindeutigen Minimierer $u_\star \in \mathcal{A}$.
\end{lemma}

\begin{proof}
Sei $\{u_n\}$ eine minimalisierende Folge. Wegen (K) ist sie beschränkt, also
besitzt sie eine schwach konvergente Teilfolge $u_{n_k} \rightharpoonup u_\star$
mit $u_\star \in \mathcal{A}$ aufgrund der schwachen Abgeschlossenheit konvexer
Mengen. Konvexität (C) plus starke Unterhalbstetigkeit (L) implizieren zusammen
die schwache Unterhalbstetigkeit von $J$ (Satz von Mazur; siehe
\cite[Prop.~II.1.2]{EkelandTemam1999}), weshalb
$J(u_\star) \leq \liminf J(u_{n_k}) = \inf_{\mathcal{A}} J$, und
$u_\star$ ist ein Minimierer. Strikte Konvexität (C) schließt einen zweiten
Minimierer aus: wenn $v_\star \neq u_\star$ ebenfalls das Minimum realisierte,
würde der Mittelpunkt $\tfrac{1}{2}(u_\star + v_\star) \in \mathcal{A}$
$J$ strikt verringern, was der Minimalität widerspricht.
\end{proof}

Der Inhalt von Lemma~\ref{lem:ucu} ist klassisch, doch die Tragweite seiner
Rolle im vorliegenden Programm ist es nicht: Es ist dasselbe Lemma,
mit verschiedenen $(\mathcal{H}, \mathcal{A}, J)$, das in jeder der
unten aufgelisteten fünf Realisierungen den eindeutigen Sättigungszustand liefert.

\subsection{Die parallelen Realisierungen}
\label{sec:ucu-instances}

Tabelle~\ref{tab:ucu-instances} führt die fünf Instanzen auf.
Jede Spalte zeigt, welches Objekt die Rolle des Hilbertraums $\mathcal{H}$,
der zulässigen Klasse $\mathcal{A}$ und des Funktionals $J$ spielt. Die letzte Zeile
vermerkt den eindeutigen Minimierer $u_\star$ und seine physikalische oder arithmetische
Interpretation. Die Tabelle ist ein Template-Ledger: Sie behauptet nicht, dass
jede Zeile denselben Beweisstatus besitzt, sondern macht sichtbar, welche
UCU-Rolle in welcher externen Argumentarchitektur verwendet wird.

\renewcommand{\arraystretch}{1.25}
\begin{table}[htbp]
\centering
\scriptsize
\setlength{\tabcolsep}{3pt}
\begin{tabular}{|p{0.16\linewidth}|p{0.40\linewidth}|p{0.34\linewidth}|}
\hline
\textbf{Programm} & \textbf{Funktional $J$ auf Klasse $\mathcal{A}$} &
\textbf{Eindeutiger Minimierer $u_\star$} \\ \hline

\emph{Riemann Even Dominance} \cite{Geiger2026RHII}
& $J(f) = \langle f, Q_\zeta f \rangle / \|f\|^2$ auf
$\mathcal{A} = L^2_\mathrm{even}([-L,L]) \setminus \{0\}$;
Konvexität auf geraden Testfunktionen reflektiert die Positiv-Definitheit
der Weilschen quadratischen Form nach der Direkten Frontier-Dominanz.
& Der Grundzustand $f_\star$, der die Spektrallücke (spectral gap) realisiert;
seine Existenz und Eindeutigkeit sind genau die
Gerade-Ungerade-Dominanzaussage. \\ \hline

\emph{Selberg Zeta (Theorem~\ref{thm:selberg-hp})}
& $J(\phi) = \int_{\Gamma\backslash\mathbb{H}} |\nabla \phi|^2$
auf $\mathcal{A} = \{\phi \in H^1_0(\mathcal{F}) : \|\phi\|_{L^2}=1\}$;
der Rayleigh-Quotient ist auf der Sphäre nach Herausdividieren
der $S^1$-Phase strikt konvex.
& Die erste Laplace--Beltrami-Eigenfunktion $\phi_1$, deren
Casimir-Eigenwert den spektralen Hilbert--P\'olya-Inhalt liefert. \\ \hline

\emph{Turbulenz / Kolmogorov K41}
& $J(u) = \int |\nabla u|^2 - \epsilon \int u$ auf
$\mathcal{A} = $ zulässige divergenzfreie Geschwindigkeitsfelder mit
fixierter Energiedissipation $\epsilon$; strikte Konvexität gilt im
Inertialbereich via Exponent-$2/3$-Skalierung.
& Das eindeutige Skalierungsprofil, das die $k^{-5/3}$-Energie-Kaskade
erzeugt \cite{FSTTurbulence}. \\ \hline

\emph{Yang--Mills Gibbs-Maß}
& $J(A) = S_{\mathrm{YM}}(A) = \tfrac{1}{4g^2}\int F \wedge {\star F}$
auf $\mathcal{A} = \mathcal{A}_{\mathrm{flat}} / \mathcal{G}$, dem
Modulraum eichäquivalenter Konnektionen; strikte Konvexität auf
dem Quotienten folgt aus der Coulomb-Eichbedingung.
& Der Vakuumsektor, der die Massenlücke
erreicht \cite{FSTYangMills}. \\ \hline

\emph{CRM (Flow-Zeta, \S\ref{sec:crm-flow})}
& $J(\mathrm{Rapidity}) = V_{\mathrm{eff}}(\mathrm{Rapidity})$, ein
effektives Dunkle-Energie-Potenzial auf $(-1,1)$ mit Rand-Koerzivität für $|x| \to 1$,
die durch die Lorentz-Boost-Komposition $\oplus$ eingebaut ist; strikte Konvexität
ist der strukturelle Input der Sättigungsaxiome \cite{CRMV}.
& Der eindeutige saturierte Boost-Zustand, der das beobachtete $\tanh$-Profil in
$w(z)$ erzeugt. \\ \hline

\end{tabular}
\caption{UCU-Instanzen als Template-Ledger. Dasselbe Lemma, verschiedene
$(\mathcal{H},\mathcal{A},J)$; der Beweisstatus der Zeile wird vom jeweiligen
externen Programm geerbt.}
\label{tab:ucu-instances}
\end{table}
\renewcommand{\arraystretch}{1.0}

\subsection{UCU und Pattern A / RFEP: Die Übersetzung}
\label{sec:ucu-pattern-a}

Die physikalische Seite von UCU im Renormierten Prinzip der freien Energie (RFEP)
\cite{RFEPFramework} ist \emph{Pattern~A}: funktionale Positivität unter einer
Eichbedingung (Gauge Constraint). Die folgende Liste ist das formale
Übersetzungstemplate:

\begin{itemize}
\item \textbf{Pattern~A (PA1).} Ein Zielfunktional $J$ (freie Energie,
Weil-Form, Rayleigh-Quotient, K41-Wirkung, Yang--Mills-Wirkung) ist
positiv-definit auf einem eich-eingeschränkten Zustandsraum.\\
$\longleftrightarrow$ \emph{UCU Hypothese (C).}

\item \textbf{Pattern~A (PA2).} Das Funktional ist nach unten beschränkt
und wächst im Unendlichen der zulässigen Klasse an.\\
$\longleftrightarrow$ \emph{UCU Hypothese (K).}

\item \textbf{Pattern~A (PA3).} Es existiert eine eindeutige stabile Sättigung
und diese wird angenommen.\\
$\longleftrightarrow$ \emph{UCU Schlussfolgerung.}
\end{itemize}

\begin{remark}[UCU vs.\ DS1--DS3]\label{rem:ucu-vs-ds}
Das Dissipative-Selektions-Prinzip DS1--DS3 von
\cite{RFEPFramework} ist der \emph{dynamische} Partner von UCU: UCU
identifiziert den eindeutigen Minimierer; DS1--DS3 identifiziert die eindeutige
Trajektorie dorthin. In allen fünf Instanzen aus
Tabelle~\ref{tab:ucu-instances} erlaubt eine zeitabhängige oder flussparametrisierte
Version von $J$ die DS1--DS3-Struktur, doch der Minimierer
selbst ist bereits durch UCU allein festgelegt. DS1--DS3 sind daher
logisch nachgelagert zu UCU.
\end{remark}

\subsection{UCU und die Trinität}
\label{sec:ucu-trinity}

Mit etabliertem UCU lautet die Trinitätsstruktur dieses Papers:
\[
  \underbrace{\text{UCU}}_{\text{Eindeutigkeit der Sättigung}}
  \;+\;
  \underbrace{\text{SGE}}_{\text{HP-Status-Klassifikation}}
  \;+\;
  \underbrace{\text{Weil-Form Transversalität}}_{\text{zweig-unabhängiges Positivitäts-Werkzeug}}
\]

UCU liefert die Ebene \emph{"warum es eine eindeutige Antwort gibt"}; SGE
(Abschnitt~\ref{sec:sge-conjecture}) liefert die Ebene \emph{"zu welchem Zweig gehört
die Familie"}; die Weil-Form-Architektur (\S\ref{sec:weil-transversal}) liefert die Ebene
\emph{"wie berechnet man Positivität in jedem Zweig"}. Die drei Ebenen sind
unabhängig voneinander: UCU prognostiziert auf dem HP-Spektrum weder YES noch NO;
SGE prognostiziert keine Spektrallücken; die Weil-Form-Transversalität liefert
für sich genommen keine Eindeutigkeit. Sie setzen sich zum vollständigen Programm zusammen.

\begin{remark}[Scope von UCU]\label{rem:ucu-scope}
UCU ist \emph{notwendig, aber nicht hinreichend} für eine Hilbert--P\'olya-Realisierung.
Ein eindeutiger Minimierer im Sinne von Lemma~\ref{lem:ucu} liefert nicht automatisch
einen hermiteschen Operator, dessen Spektrum die Nullstellen von $Z_X$ kodiert;
er liefert einen Grundzustand. Der Übergang vom Grundzustand zum Spektrum
erfordert die SGE-Klassifikation (um zu identifizieren, ob man eine Gruppe oder
nur eine Halbgruppe von Transferoperatoren hat) und die Weil-Form-Architektur
(um die explizite Formel-Positivität wiederherzustellen). UCU ist
der gemeinsame Nenner, nicht der Beweis von Hilbert--P\'olya in irgendeiner
einzelnen Instanz.
\end{remark}

% ============================================================
\section{Weil-Form-Methoden als branch-transversales Komplement}
\label{sec:weil-transversal}

Wir argumentieren nun, dass die Weilsche quadratische Form $Q_X$ aus
Abschnitt~\ref{sec:weil-qf} einheitlich in allen drei Zweigregimes gilt.

\subsection{Form-Definitionsbereich Verankerung}
\label{sec:weil-form-domain}

Für jede Zeta-Typ-Familie $X$ ist die Weilsche quadratische Form $Q_X$
durch das Quadrupel
\begin{equation}\label{eq:weil-quadruple}
  \bigl(Q_X,\; \operatorname{Dom}(Q_X),\; P_X,\;
        \operatorname{Dom}(Q_X)^{\pm}\bigr)
\end{equation}
verankert, wobei $\operatorname{Dom}(Q_X) \subset L^2(\mathbb{R})$ ein
abgeschlossener Teilraum ist, auf dem $Q_X$ wohldefiniert und nach unten beschränkt ist
(der Form-Definitionsbereich), $P_X$ der Paritätsoperator $f(t) \mapsto f(-t)$ ist und
$\operatorname{Dom}(Q_X)^{\pm} = \operatorname{Dom}(Q_X) \cap L^2_{\mathrm{even/odd}}$
die paritäts-invarianten Teilräume sind. Für die arithmetischen Instanzen
dieses Papers (Abschnitte~\ref{sec:riemann}--\ref{sec:primehub})
wird der Form-Definitionsbereich als der Friedrichs-Abschluss von
$C_c^{\infty}([-\log\lambda,\log\lambda])$ bezüglich der
Graphennorm von $Q_X$ angenommen, und die Kosinus-/Sinusbasen von
Abschnitt~\ref{sec:weil-qf} realisieren ein Dichtheitssystem in
$\operatorname{Dom}(Q_X)^{\pm}$. Die Paritätserhaltung
$P_X \operatorname{Dom}(Q_X) = \operatorname{Dom}(Q_X)$ ist eine
Konsequenz aus der Paritätsinvarianz des archimedischen Kerns
$e^{u/2}/|2\sinh u|$ (der Absolutbetrag wird genommen; die Weil-Formel verwendet
die Hauptwert-Konvention (principal value) bei $u = 0$, vgl. \cite{Weil1952})
und der symmetrischen Primverschiebungs-Operatorfamilie
$S_{m\log p} + S_{-m\log p}$, die in die Weilsche Form eingeht.

\begin{remark}[Scope dieser Verankerung]\label{rem:form-domain-scope}
Das explizite Form-Definitionsbereich-Lemma, das der Riemann-Instanz
zugrunde liegt (Einschränkung von $Q_\zeta$ auf einen $H^{1/2}$-Abschluss, auf dem der
archimedische Kern form-beschränkt ist, und Verifikation der Paritätserhaltung
auf dem vollen Form-Definitionsbereich), ist als Minor-Revision-Punkt
in \cite{RHCritiqueReview2026} (K5) aufgeführt. Der branch-transversale Inhalt
des vorliegenden Papers ist unabhängig von dieser Verfeinerung: Das Shift-Parity-Lemma
und der Beweis der Direkten Frontier-Dominanz operieren auf
endlich-dimensionalen Galerkin-Blöcken, und ihr Transfer auf den kontinuierlichen Operator
nutzt Cauchy-Interlacing plus Schur-Komplement-Schranken (siehe
Abschnitt~\ref{sec:three-component}, Komponente $C^{(2)}$), und nicht
die Form-Definitionsbereichs-Regularität per se.
\end{remark}

\subsection{Positivität bei HP-BL-YES}

Für $X = Z_\Gamma$ (Selberg, Theorem~\ref{thm:selberg-hp}) ist die
Positivität von $Q_{Z_\Gamma}$ auf der relevanten Testfunktionsklasse eine
Konsequenz aus der Selbstadjungiertheit von $\Delta_\mathcal{F}$; die Positivität
der Weilschen Form lässt sich aus der spektralen Zerlegung
des Laplace-Operators ableiten. Dies ist eine Konsistenzprüfung: Beide Architekturen
liefern dieselbe Antwort.

\subsection{Positivity bei HP-BL-NO}

Für $X = \zeta$ (Riemann, Theorem~\ref{thm:boundary} Zeile 1) wird die
Positivität von $Q_\zeta$ in \cite{Geiger2026RHII}
über die Direct-Frontier- / Even-Dominance-Route verfolgt
(Shift-Parity-Lemma + Mertens-Theorem + CAP-zertifizierter Basisfall).
Der aktuelle öffentliche Status ist bedingt, aber die dortige Route
nutzt keinerlei Hilbert-P\'olya-Input. Dies ist der strikt
branch-transversale Inhalt.

\subsection{Positivität bei HP-BL-OPEN}

Für $X = G_\mathrm{prime}$ zerfällt die Frage nach der Positivität der Weil-Form
in zwei Schritte. Erstens muss die explizite Formel $\mathcal{E}_\mathrm{prime}$
für den Prime-Hub mit der expliziten Formel von Riemann--Weil
\eqref{eq:weil-riemann} auf der Prim-/Geodätenseite übereinstimmen, da
die Bedingungen (OP5.1)--(OP5.3) aus Abschnitt~\ref{sec:primehub}
vorschreiben, dass die Fredholm-Determinante von $\mathcal{L}_s$ gleich $1/\zeta(s)$ ist.
Nach Proposition~\ref{prop:fredholm-orientation} muss diese
Determinantenidentität durch einen Quotienten-, regularisierten oder
gradierten Kanal interpretiert werden, wenn OP5.5 sinnvoll bleiben
soll. Sobald diese Kanalsemantik festgelegt ist, gilt
$Q_{G_\mathrm{prime}} = Q_\zeta$ als quadratische Form auf der
Testfunktionsklasse, und ihre Positivität reduziert sich auf die von Riemann
(das HP-BL-NO-Regime). Zweitens ist die Weil-Form-Positivität --- \emph{vorausgesetzt}, dass schließlich eine explizite
Prime-Hub-Konstruktion hervorgebracht wird ---
automatisch durch das branch-transversale Argument aus
\cite{Geiger2026RHII} auf der gemeinsamen Prim-Seite gegeben. Dies ist die stärkste
mögliche Ausprägung von Branch-Transversalität: Die Positivität der
quadratischen Form ist invariant unter dem $\HBL$-Statusübergang.

Naheliegende Evidenz für eine verwandte HP-BL-OPEN Familie: Die numerischen
Arbeiten aus Session~8 zu Dirichlet $L$-Funktionen
\cite{PaleyWienerDirichlet} zeigen, dass Versuche, $H_\chi$ aus phänomenologischen Kernen
(Gauss- oder Hadamard-gewichteter Pol) zu konstruieren,
systematisch scheitern, während die Positivität der Weilschen Form
zugänglich bleibt. Dies ist eine zweite empirische Instanz, in der
der HP-BL-Status faktisch $\mathrm{NO}$ oder $\mathrm{OPEN}$ ist, Weil-Form-Methoden
jedoch weiterhin anwendbar sind.

\subsection{Zusammenfassung}

Die Weilsche quadratische Form ist branch-transversal: Ihre Existenz und ihr
Positivitäts-Inhalt hängen nicht vom HP-BL-Status der Familie ab.

% ============================================================
\section{Die Halbgruppen-Gruppen-Äquivalenz (SGE)}
\label{sec:sge-conjecture}

\subsection{Eine vereinheitlichende Vermutung}
\label{sec:sge-conjecture-statement}

Das Randtheorem~\ref{thm:boundary} zeigt drei qualitativ distinkte Werte
$\{\mathrm{YES}, \mathrm{NO}, \mathrm{OPEN}\}$ von $\HBL$ an drei kanonischen Instanzen
auf. Eine natürliche Frage ist, ob diese Trichotomie \emph{fundamental}
oder bloß eine Oberflächeneigenschaft eines tieferen einzigen Kriteriums ist.
Wir formulieren die folgende vereinheitlichende Vermutung.

Sei $(Z_X, \mathcal{T}_X, \mathcal{E}_X)$ eine Zeta-Typ-Familie im Sinne von
Definition~\ref{def:zeta-family}. Die Transferfamilie $\mathcal{T}_X$ wird
typischerweise durch eine diskrete Menge $I_X$ von „Primobjekten“ organisiert
(rationale Primzahlen, geschlossene Geodäten, primitive Zyklen), von denen jedes
eine kanonische Länge $\ell(\iota)$ trägt. Die zugehörige Transfer-Provenienz
liefert entweder ein halbgruppenartiges Objekt (Multiplikation von
Primzahlpotenzen, Idealmonoid, Verkettung primitiver Zyklen) oder eine
Gruppenwirkung mit unitärer Darstellung und Spurformel. In der Selberg-Zeile
ist die Gruppenseite kein freies Kompositionsgesetz auf geschlossenen
Geodäten; geschlossene Geodäten gehen als Spurformel-Daten ein, während der
Casimir aus der $\PSL_2(\mathbb{R})$-invarianten Geometrie stammt.

\begin{conjecture}[Halbgruppen-Gruppen-Äquivalenz, SGE]
\label{conj:sge}
Für die in diesem Paper kalibrierte Klasse von Zeta-Typ-Familien gilt:
\[
  \HBL(X) = \mathrm{YES}
  \quad\Longleftrightarrow\quad
  (I_X, \cdot)\ \text{hat lokalkompakte Gruppen-Provenienz}.
\]
\end{conjecture}

Die Rückrichtung (Gruppe $\Rightarrow$ YES) ist nur in den realisierten
Spurformel-Fällen klassisch: Eine natürliche Gruppenwirkung zusammen mit
der passenden unitären Darstellung und spektralen Spurformel produziert
einen Casimir-artigen Operator, der (HP1)--(HP3) erfüllt. Die
SGE-Vermutung ist die stärkere Behauptung, dass dieser
Modellfall-Mechanismus die branch-lokale Hilbert-P\'olya-Grenze exakt
erfasst. Die Hinrichtung (Halbgruppe $\Rightarrow$ NO) ist der
substanzielle Inhalt von Vermutung~\ref{conj:sge} und steht im Einklang
mit jeder bisher getesteten Instanz:

\begin{center}
\begin{tabular}{@{}>{\raggedright\arraybackslash}p{0.24\textwidth}
                  >{\raggedright\arraybackslash}p{0.42\textwidth}
                  >{\raggedright\arraybackslash}p{0.25\textwidth}@{}}
\toprule
Familie & Transfer-Algebra $(I_X, \cdot)$ & Prognostiziertes $\HBL$ \\
\midrule
Riemann $\zeta(s)$ & $\{p^m\}$ (kommutative Halbgruppe) & NO$_{\mathrm{prime}}^{\mathrm{cond}}$ (Thm.~\ref{thm:boundary}) \\
Dirichlet $L(s,\chi)$ & $\{p^m : p \nmid q\}$ (Unter-Halbgruppe) & SGE-prognostiziertes NO; kein einfacher perturbativer HP-Operator in \cite{PaleyWienerDirichlet} \\
Dedekind $\zeta_K(s)$ & Primideale $\{\mathfrak{p}\}$ (Ideal-Monoid) & SGE-prognostiziert NO; endliche Evidenz \\
Selberg $Z_\Gamma(s)$ & $\PSL_2(\mathbb{R})$-Geometrie; Geodäten als Spurdaten & YES (Thm.~\ref{thm:selberg-hp}) \\
Ihara $\zeta_G(u)$ & geschlossene Wege in ungerichtetem $G$ & endliches YES-Zertifikat/Proxy (Petersen) \\
Automorph $L(s,\pi)$ & Hecke-Algebra auf $\Gamma\backslash G$ & Kandidaten-YES (spurformelabhängig) \\
Prime-Hub $G_\mathrm{prime}$ & primitive Zyklen (freie Halbgruppe) & SGE-prognostiziert NO; aktueller Status OPEN \\
\bottomrule
\end{tabular}
\end{center}

Trifft Vermutung~\ref{conj:sge} zu, kollabiert die scheinbare Trichotomie
von Theorem~\ref{thm:boundary} zu einem einzigen algebraischen Kriterium.
Der OPEN-Status von Prime-Hub ist in dieser Lesart kein dritter
fundamentaler Modus, sondern ein Zustand unvollständigen Wissens: die explizite
Konstruktion von $(I_\mathrm{prime}, \cdot)$ würde deterministisch zwischen
YES und NO via Vermutung~\ref{conj:sge} entscheiden. Eine separate
Arbeitsnotiz \cite{SGENote} entwickelt die präzise Formulierung und
sammelt Falsifikationskriterien.

\subsection[Erster empirischer Test von Vermutung: Dedekind-Zeta von Q(sqrt(-5))]{Erster empirischer Test von Vermutung~\ref{conj:sge}:
Dedekind-Zeta von \texorpdfstring{$\mathbb{Q}(\sqrt{-5})$}{Q(sqrt(-5))}}
\label{sec:sge-dedekind-test}

Als direkten Test von Vermutung~\ref{conj:sge} außerhalb der drei kanonischen
Instanzen aus Theorem~\ref{thm:boundary} haben wir den folgenden Test
\cite{DedekindNEBTest} durchgeführt. Für den imaginär-quadratischen Körper
$K = \mathbb{Q}(\sqrt{-5})$ (Diskriminante $-20$, Klassenzahl $2$) hat die
Dedekind-Zetafunktion $\zeta_K$ eine Transfer-Algebra, die über
Primideale $\mathfrak{p}$ von $\mathcal{O}_K = \mathbb{Z}[\sqrt{-5}]$ indiziert
wird; diese bildet unter Idealmultiplikation eine Halbgruppe.
Vermutung~\ref{conj:sge} prognostiziert daher $\HBL(\zeta_K) = \mathrm{NO}$.

Um dies zu testen, enumerierten wir die $92$ Primideale von $K$ mit Norm bis
höchstens $500$, klassifiziert in $2$ verzweigte (über rationalen Primzahlen $2, 5$),
$86$ zerlegte (zwei Ideale für jede rationale Primzahl $p \equiv 1, 3, 7,
9 \pmod{20}$) und $4$ träge (Norm $p^2$ für $p \equiv 11, 13, 17,
19 \pmod{20}$). Die entsprechenden normalisierten Verschiebungen $r_\mathfrak{p}
= 2\log N\mathfrak{p}/L$ mit $L = 10$ ergaben $49$ distinkte Werte
in $(0, 2)$ --- weniger als die $95$, die für die rationalen Primzahlen
von $\mathbb{Q}$ bis zur selben Normschranke erzielt wurden, bedingt durch die Verdopplung
zerlegter Primzahlen, wodurch Idealpaare auf einen einzelnen Verschiebungswert kollabierten.
Für jede Verschiebung konstruierten wir die NE-B-Matrix $D_N(r)$ aus Definition~\ref{thm:NE-B}
mittels numerischer Quadratur, und berechneten die symmetrische gemeinsame
Zentralisator-Dimension, indem wir das lineare System
$[M, D_N(r_\mathfrak{p})] = 0$ für die $N(N+1)/2$ symmetrischen
Freiheitsgrade lösten.

\begin{center}
\begin{tabular}{@{}c c c c@{}}
\toprule
$N$ & $\mathbb{Q}$ (Kontrolle) & $\mathbb{Q}(\sqrt{-5})$ & zufällige Verschiebungen \\
\midrule
$5$ & $\dim(Z) = 1$ & $\dim(Z) = 1$ & $\dim(Z) = 1$ \\
$7$ & $\dim(Z) = 1$ & $\dim(Z) = 1$ & $\dim(Z) = 1$ \\
$10$ & $\dim(Z) = 1$ & $\dim(Z) = 1$ & $\dim(Z) = 1$ \\
\bottomrule
\end{tabular}
\end{center}

Alle drei Spalten stimmen überein: Nur der skalare Zentralisator überlebt bei
jedem getesteten $N \in \{5, 7, 10\}$. Dies stellt endliche kontrollierte Evidenz
für Vermutung~\ref{conj:sge} außerhalb der drei kanonischen Instanzen von
Theorem~\ref{thm:boundary} dar: Trotz des qualitativ distinkten
Primzerlegungsmusters von $K$ und der grob um den Faktor zwei niedrigeren Verschiebungsdichte relativ zu $\mathbb{Q}$,
weist die Dedekindsche Transfer-Algebra von $\mathbb{Q}(\sqrt{-5})$
dasselbe NE-B-Verhalten wie die Familie rationaler Primverschiebungen auf.
Die Beobachtung stützt die algebraische (anstatt dichte-basierte) Natur von NE-B
und, transitiv, der Prognose $\HBL = \mathrm{NO}$.

\subsection{Komplementärer empirischer Test von
Vermutung~\ref{conj:sge}: Ihara-Zeta des Petersen-Graphen}
\label{sec:sge-ihara-test}

Auf der $\HBL = \mathrm{YES}$-Seite der SGE-Vermutung wird ein direkter
Test durch die Ihara-Zetafunktion $\zeta_G(u)$ eines ungerichteten
endlichen Graphen $G$ bereitgestellt. Die Transfer-Algebra in diesem Fall
ist das Wege-Gruppoid modulo Backtracking (Umkehrung), dessen assoziierte Struktur die
Fundamentalgruppe $\pi_1(G)$ ist --- eine freie Gruppe $F_r$ mit $r = m - n + 1$
Erzeugern. Da $F_r$ eine Gruppe (nicht bloß eine Halbgruppe) ist,
prognostiziert Vermutung~\ref{conj:sge} einen YES-seitigen Mechanismus
für Graph-Zetafunktionen. Der endliche Test unten zertifiziert nur die
Petersen-Instanz als kontrollierten Proxy; er ist kein Theorem für jeden
zusammenhängenden ungerichteten Graphen.

Wir testeten diese Prognose am Petersen-Graphen (3-regulär, $n = 10$, $m = 15$,
ein Ramanujan-Graph mit $\mathrm{Aut}(G) = S_5$). Drei unabhängige
Verifikationen wurden durchgeführt \cite{IharaPetersenTest}:
\begin{enumerate}
\item[(V1)] \textbf{Bass-Hashimoto Identität.} Die klassische Identität
  $\det(I - uB) = (1 - u^2)^{m-n}\det(I - Au + Qu^2)$, wobei $B$ der
  Hashimoto-Non-Backtracking-Operator auf gerichteten Kanten ist und
  $Q = D - I$, wurde am Testpunkt $u = 0.3 + 0.1 i$ bis
  auf 15 signifikante Ziffern verifiziert.
\item[(V2)] \textbf{Ramanujan RH-Analogon.} $18$ der $19$ nicht-trivialen Nullstellen
  von $\zeta_G(u)^{-1}$ liegen exakt auf dem kritischen Kreis
  $|u| = 1/\sqrt{q-1} = 1/\sqrt{2}$. Die einzige Ausnahme bei $|u| = 1/2$
  entspricht dem trivialen Adjazenz-Eigenwert $\lambda_A = 3$.
\item[(V3)] \textbf{Endliches YES-Zertifikat/Proxy.} Der Graph-Laplace-Operator
  $\Delta_G = D - A$ ist selbstadjungiert mit Spektrum
  $\{0, 2^{(5)}, 5^{(4)}\}$. Seine Eigenwerte hängen mit den Ihara-Nullstellen
  über das quadratische Polynom $2u^2 - \lambda_A u + 1 = 0$ zusammen.
  Jeder Graph-Automorphismus $\sigma \in S_5$ induziert eine Permutationsmatrix $P_\sigma$,
  die mit $\Delta_G$ kommutiert, und der Kanten-Lift $P_\sigma^{\mathrm{edge}}$ kommutiert mit $B$.
  Der Zentralisator von $B$ in $\mathrm{End}(\mathbb{C}^{2m})$ enthält daher
  die $S_5$-Darstellungsalgebra, deren nicht-skalare Elemente ein
  endliches YES-Zertifikat/Proxy für
  $\HBL(\zeta_{\mathrm{Petersen}})$ liefern.
\end{enumerate}

In Kombination mit dem Dedekind-Test aus Abschnitt~\ref{sec:sge-dedekind-test}
wurde Vermutung~\ref{conj:sge} nun mit endlichen Kontrollen auf \emph{beiden}
Seiten der Halbgruppen-Gruppen-Dichotomie stressgetestet: auf der Halbgruppenseite
($\mathbb{Q}(\sqrt{-5})$, $\dim(Z) = 1$) bei $N \leq 10$, und auf der
Gruppenseite (Petersen, nicht-skalare kommutierende Operatoren explizit
konstruiert). Fünf qualitativ distinkte Zeta-Typ-Familien
stützen nun konsistent Vermutung~\ref{conj:sge} (drei kanonische Instanzen plus
dieser beiden Tests), und es wurde kein Gegenbeispiel erbracht. Dies bleibt
endliche Evidenz, kein Beweis.
Die Anhebung zu einem formal bewiesenen Theorem würde ein strukturelles
Argument (anstatt weiterer Fallstudien) erfordern und ist als primärer nächster
Schritt identifiziert.

\begin{remark}[Diskriminierendes Kontrollexperiment]
\label{rem:sge-control}
Eine berechtigte Sorge bezüglich des Dedekind-Tests aus
Abschnitt~\ref{sec:sge-dedekind-test} ist, dass $\dim(Z) = 1$ die
\emph{generische} Zentralisator-Dimension von beliebigen Matrixfamilien sein
könnte und nicht eine strukturelle Signatur von SGE-NO. Ein dediziertes
Kontrollexperiment \cite{SGEControlExperiment2026} adressiert dies,
indem dieselbe Zentralisator-Dimensions-Routine an expliziten SGE-YES
Familien (zyklische Gruppe $\mathbb{Z}/N$ reguläre Darstellung;
elementar-abelsche $(\mathbb{Z}/2)^k$) neben
einer zufällig-symmetrischen Baseline laufen lässt. Das Ergebnis ist eindeutig:
\begin{itemize}
\item $\mathbb{Z}/N$: $\dim(Z_{\text{sym}}) = 7, 10, 14, 17, 22$
  bei $N = 5, 7, 10, 12, 15$ (wächst linear mit $N$).
\item $(\mathbb{Z}/2)^k$: $\dim(Z_{\text{full}}) = 2^k$ exakt
  (gleich der Gruppenordnung) für $k = 2, 3, 4$.
\item Zufällig-symmetrische Baseline: $\dim(Z_{\text{sym}}) = 1$ für
  alle getesteten $N$.
\end{itemize}
Der Testapparat diskriminiert folglich sauber zwischen SGE-YES
(Zentralisator-Dimension skaliert mit $N$) und SGE-NO / Null-Fall
(Dimension eins). Das Dedekind-Ergebnis $\dim(Z) = 1$ ist mithin
strukturell informativ und kein generisches Artefakt. Dies schließt eine
Dokumentationslücke, die im 7-Phasen-Review \cite{RHCritiqueReview2026}
dieses Manuskripts identifiziert wurde.
\end{remark}

Connes' adelisches Programm \cite{Connes1999, Connes2026} passt
natürlich in dieses Bild: Indem es die Halbgruppe $\{p^m\}$ durch die
Ideleklassengruppe $\mathbb{A}_K/K^\times$ ersetzt, \emph{transportiert} es die
Riemann-Familie aus der Halbgruppenspalte in die Gruppenspalte der
Tabelle. In der SGE-Lesart würde Connes' Programm
$\HBL(\zeta)$ von NO (in der klassischen Primverschiebung-Transfer-Algebra)
zu YES (in der adelischen Transfer-Algebra) anheben, ohne
Theorem~\ref{thm:boundary} zu widersprechen.

% ============================================================
\section{Die Blockchain-Lesart: eine spekulative Brücke}
\label{sec:blockchain}

\emph{Dieser Abschnitt ist als eine \textbf{spekulative Brücke} in der
Vier-Klassen-Rigoroshierarchie von Abschnitt~\ref{sec:rigor} markiert: Er
überträgt den mathematischen Inhalt der Trinität in ein informatisches Register
durch eine Korrespondenz mit dem Bitcoin-Blockchain-Protokoll. Keine
mathematische Behauptung in diesem Abschnitt ist neu; der Beitrag ist eine
konsolidierte Lesart, die die Architektur der Trinität
Lesern zugänglich macht, die mit kryptografischen Protokollen vertraut sind.}

\subsection{Zwei algebraische Welten in einem Protokoll}
\label{sec:blockchain-two-worlds}

Das Bitcoin-Protokoll nutzt simultan zwei algebraische Strukturen,
die aus der Perspektive von Vermutung~\ref{conj:sge} von
entgegengesetztem Typ sind:

\begin{enumerate}
\item Die \textbf{Hash-Funktion} (SHA-256) generiert eine Bildmenge
  mit einer \emph{Halbgruppenstruktur} unter der Verkettung von
  Datenblöcken. Der Hash ist absichtlich eine Einwegfunktion --- er hat keine
  natürliche inverse Operation --- also ist die Indexmenge einer Hash-Kette
  eine Halbgruppe, keine Gruppe.
\item Die \textbf{digitale Signatur} (ECDSA auf der elliptischen Kurve
  $\mathrm{secp256k1}$) operiert auf einer \emph{abelschen Gruppe}: Die
  $E(\mathbb{F}_p)$-Punkte bilden eine Gruppe unter Tangenten-und-Sekanten-Addition (chord-and-tangent addition),
  mit expliziten Inversen und einer wohldefinierten Casimir-ähnlichen
  Invariante (die Kurvengleichung $y^2 = x^3 + 7$).
\end{enumerate}

In der Sprache von Abschnitt~\ref{sec:sge-conjecture} lebt SHA-256
auf der Halbgruppenseite der SGE-Dichotomie, während ECDSA auf
der Gruppenseite lebt. Das Bitcoin-Protokoll ist somit eine \emph{physikalische
Realisierung} der algebraischen Dichotomie, die Vermutung~\ref{conj:sge}
auf der Ebene der Zeta-Typ-Familien identifiziert, wobei die Hash-Seite
die Riemannsche Primzahlpotenz-Halbgruppe $\{p^m\}$ spiegelt und die
Signatur-Seite die Selbergsche Geodäten-Gruppenstruktur
$\Gamma \subset \PSL_2(\mathbb{R})$.

\subsection{Die Korrespondenz-Tabelle}
\label{sec:blockchain-correspondence}

\begin{center}
\renewcommand{\arraystretch}{1.25}
\begin{tabular}{@{}p{0.38\linewidth} p{0.52\linewidth}@{}}
\toprule
\textbf{Baum-theoretische Struktur (Tree-theoretic structure)} & \textbf{Bitcoin-Protokoll Analogon} \\
\midrule
Primzahlpotenz-Halbgruppe
  $I_\zeta = \{p^m : p \text{ prim}, m \geq 1\}$ &
SHA-256 Hash-Raum: Halbgruppe unter Datenverkettung, kein
natürliches Inverses, liefert eindeutige Prüfsummen (analog zur eindeutigen Primfaktorzerlegung). \\[0.4ex]

Hilbert-P\'olya Gruppe
  $I_{Z_\Gamma} = \Gamma \subset \PSL_2(\mathbb{R})$ &
ECDSA Signatur-Gruppe: $E(\mathbb{F}_p)$-Punkte auf
$\mathrm{secp256k1}$ mit Tangenten-und-Sekanten-Addition, Inverse
explizit, Casimir-ähnliche Kurveninvariante $y^2 = x^3 + 7$. \\[0.4ex]

Weilsche quadratische Form $Q_X$ (branch-transversale Architektur; liefert
Lücken-Positivität (gap positivity) ungeachtet des $\HBL$-Status) &
Nakamoto-Konsens / Proof-of-Work: eine Funktion auf Protokollebene, die
die Kette einheitlich über alle Hash-seitigen und Signatur-seitigen
Teilnehmer validiert, unabhängig davon, welche algebraische Struktur
eine individuelle Transaktion nutzt. \\[0.4ex]

UCU: Strikte Konvexität des maßgeblichen Funktionals liefert einen
eindeutigen Minimierer (Abschnitt~\ref{sec:ucu}) &
Längste-Kette-Regel (longest-chain rule): strikt zunehmende Schwierigkeit (difficulty) entlang der schwersten
Kette liefert einen eindeutigen Konsens-Kopf; kumulativer Proof-of-Work ist
der strikt konvexe „Score“, der ihn auswählt. \\[0.4ex]

Theorem~\ref{thm:NE-B}: kein universeller kommutierender Operator für die
Primverschiebung-Familie an getesteten Galerkin-Trunkierungen &
Dezentralisierung: Kein einzelner Bitcoin-Knoten hält einen universellen Validierungsschlüssel;
das Kollektiv (Primzahlen / Miner) validiert durch unabhängige
Beiträge, nicht durch ein gemeinsames Geheimnis (shared secret). \\[0.4ex]

Riemannsche Vermutung: alle nicht-trivialen Nullstellen auf der kritischen Gerade &
Ketten-Konsistenz: alle akzeptierten Blöcke liegen auf der kanonischen
(längsten, nach Proof-of-Work schwersten) Kette. \\[0.4ex]

Dreikomponenten-Zerlegung
  $\mathrm{RH}_X = \mathrm{GBC}_X + C^{(2)}_X + \mathrm{Hurwitz}_X$
  \cite{Geiger2026RHII} &
Drei-Schichten-Protokoll-Stack: Transaktionsschicht (GBC-Analogon, spektraler
Inhalt des Operators), Konsens-Schicht ($C^{(2)}$-Analogon, Connes-Approximation),
historische Persistenz (Hurwitz-Analogon, Realität der
Nullstellen bewahrt die Kette über Reorganisationen hinweg). \\[0.4ex]

Vermutung~\ref{conj:sge}: $\HBL$ ist eine binäre algebraische Invariante &
Protokoll-Hybridismus: Das Funktionieren von Bitcoin erfordert die Kooperation
einer Halbgruppen-seitigen Komponente (Hashes) und einer Gruppen-seitigen
Komponente (Signaturen). Jede für sich ist unzureichend; ihre
branch-transversale Kopplung (das Protokoll selbst) ist das Analogon der
Weilschen quadratischen Form. \\
\bottomrule
\end{tabular}
\end{center}

\subsection{Die Trinität kryptografisch gelesen}
\label{sec:blockchain-trinity}

Die drei Metaprinzipien aus Abschnitt~\ref{sec:trinity} erhalten eine
kompakte kryptografische Lesart:

\begin{itemize}
\item \textbf{UCU als Proof-of-Uniqueness.}
Die strikte Konvexität des maßgeblichen Funktionals ist das Analogon auf
Protokollebene zu Bitcoins Längste-Kette-Regel. In beiden Fällen wird
die Eindeutigkeit des ausgewählten Zustands durch eine monotone Bewertungsfunktion
garantiert (Difficulty-akkumulierte Arbeit für Bitcoin; das Weil-Funktional
für die Zeta-Familie). Nicht-Eindeutigkeit würde in einem Fork resultieren; die
Konvexitätsbedingung verhindert Forks in beiden Domänen.

\item \textbf{SGE als Protokoll-Hybridismus.}
Vermutung~\ref{conj:sge} postuliert, dass jede Zeta-Familie definitiv auf
einer Seite der algebraischen Dichotomie liegt. Bitcoin hingegen ist ein einzelnes
Protokoll, das beide Seiten nutzt: seine Hash-Schicht ist $\HBL = \mathrm{NO}$-strukturiert,
seine Signatur-Schicht ist $\HBL = \mathrm{YES}$-strukturiert. In dieser Lesart ist
Bitcoin der Prototyp eines \emph{hybriden Protokolls} --- ein Objekt, dessen Gesamtarchitektur
für eine einzelne Zeta-Familie unmöglich ist, aber natürlich ist,
wenn die Halbgruppen- und Gruppen-Seiten über einen gemeinsamen
branch-transversalen Mechanismus gekoppelt sind.

\item \textbf{Weil-Form als branch-transversaler Konsens.}
Die Weilsche quadratische Form $Q_X$ funktioniert einheitlich für
$\HBL = \mathrm{YES}$ und $\HBL = \mathrm{NO}$ Familien. Der Nakamoto-Konsens
funktioniert auf ähnliche Weise einheitlich für Hash-seitige Operationen und Signatur-seitige
Operationen. In beiden Fällen ist die branch-transversale
Schicht das, was den zwei algebraischen Welten eine sinnvolle Kooperation ermöglicht.
\end{itemize}

\subsection{Warum die Korrespondenz eine kontrollierte Analogie ist}
\label{sec:blockchain-substance}

Drei Beobachtungen machen diese Lesart strukturell nützlich, ohne sie in ein
Theorem zu verwandeln:

\begin{enumerate}
\item \emph{Die Dichotomie ist präzise.} Bitcoins Hash-/Signatur-Architektur
ist eine strukturelle Notwendigkeit, keine Designentscheidung: Jedes
kryptografische Konsenssystem kombiniert typischerweise eine Einweg-Digest-Funktion
(Halbgruppe) mit einer invertierbaren Authentifizierungsfunktion
(Gruppe). Das spiegelt, beweist aber nicht,
Vermutung~\ref{conj:sge}'s binäre Klassifikation.

\item \emph{Konsens ist die Algebra.} Nakamotos Einsicht war, dass ein
Proof-of-Work-Konsens als branch-transversale Schicht fungieren kann, die
die beiden Seiten vereint. Die Weilsche quadratische Form spielt dieselbe Rolle
im Zeta-Familien-Baum: sie ist die gemeinsame Arithmetik, auf der beide
$\HBL$-Regimes operieren.

\item \emph{Das Eindeutigkeitstheorem wird geteilt.} Bitcoins Längste-Kette-Regel
ist \emph{mathematisch} ein strikt konvexes Auswahlprinzip;
die strikte Konvexität des Weil-Funktionals ist dasselbe Prinzip im
kontinuierlichen Setting. Beide sind Instanzen eines allgemeinen
Konvexität-impliziert-Eindeutigkeit-Arguments (UCU).
\end{enumerate}

Dementsprechend ist die Korrespondenz am besten als strukturelles Echo zu lesen.
Dieselbe Trinität --- UCU, SGE, Weil-Form-Transversalität --- liefert eine
Kommunikationsanalogie zwischen der Klassifikation der Zeta-Typ-Programme und
der Architektur der Bitcoin-Blockchain.

\subsection{Rigorositäts-Hinweis und weiterführende Literatur}
\label{sec:blockchain-disclaimer}

Der vorliegende Abschnitt ist eine \emph{spekulative Brücke}. Er führt kein
neues mathematisches Theorem und kein neues kryptografisches Resultat ein; die
obigen Korrespondenzen sind konzeptionelle Paarungen, keine Identitäten.
Was der Abschnitt bietet, ist ein konsolidiertes Register, in dem
der strukturelle Inhalt der Trinität Lesern
aus Kryptografie, Informatik und der Erforschung dezentraler Systeme
vermittelt werden kann. Für eine umfassendere informelle Entwicklung der Blockchain-Analogie
im breiteren Baum-der-Zeta-Familien-Kontext (inklusive
verwandter Motive über Primzahlen als holografische Randdaten und über Zeit
als Informationsakkumulation), siehe die Begleitnotizen
\cite{MetaGalerkinBridge}, Appendix~I8--I10, und die dort referenzierten ergänzenden Arbeitsnotizen.

% ============================================================
\section{Diskussion und Ausblick}
\label{sec:outlook}

\subsection{Konsequenzen für Dirichlet \texorpdfstring{$L$}{L}-Funktionen}

Arbeiten aus Session~8 zu Dirichlet $L$-Funktionen \cite{PaleyWienerDirichlet}
haben etabliert, dass ein Paley-Wiener Transfer mittels phänomenologischer Kerne
(Gauss- oder Hadamard-gewichteter Pol) die
charakter-spezifischen Lücken-Signaturen nicht erfasst. Das Randtheorem
(Theorem~\ref{thm:boundary}) interpretiert dieses Resultat: Für reelle Dirichlet-Charaktere
mit kleinem Führer (low-conductor) wird kein natürlicher HP-BL Operator
durch einfache Perturbation des Riemannschen Falls konstruiert, und die beobachteten
charakter-spezifischen Lücken müssen aus der branch-transversalen
Weil-Form-Architektur entstehen, deren Reduktionen führender Ordnung nachweislich
scheitern. Ein dediziertes Dirichlet-L Atlas Paper
\cite{DirichletAtlas} dokumentiert dieses Scheitern im Detail.

Unter Vermutung~\ref{conj:sge} ist dieses Ergebnis strukturell
prognostiziert: Die Dirichletsche Transfer-Algebra ist eine Halbgruppe, folglich
gilt $\HBL(L(\cdot, \chi)) = \mathrm{NO}$ und es kann kein natürlicher HP-Operator
existieren. Der Dirichlet-Atlas ist daher kein Fehlschlag des
Programms, sondern eine Bestätigung seiner architektonischen Grenzen.

Eine Post-Zookeeper-Verfeinerung ist nun Teil des CORE-Status. Der
Spectral Zookeeper \cite{Zookeeper2026} liefert ein unabhängiges Kern-Paket
für die Riemann-Zelle: Rang-eins-Zerlegung (rank-one decomposition), Cauchy
Interlacing, Drei-Lemma-Abschluss (three-lemma closure) und spektrale Mikrocluster. Dies
ändert nicht das Randtheorem; es verfeinert die UBA-Ebene. Die
Riemann-Zelle trägt nun zwei orthogonale Kern-Pakete: v2.1 Even
Dominance als die publizierte C2-sensitive Schicht, und das Zookeeper
Drei-Lemma-Paket als der Abschluss des Riemann-Schritts $C^{(2)}$/MS2
auf Theorem-Ebene im CCM/Fourier-Zweig. Für die Dirichlet-Zelle
ist der natürliche Rücktransfer die Route D:
ein $\chi$-getwisteter CCM-Operator mit einem Führer-Rang-eins-Term (conductor rank-one term) und einer
$\chi$-gewichteten arithmetischen Perturbation. Folglich sollte der Dirichlet-Atlas
als der erste konkrete UBA-Stresstest dieses Transfers gelesen werden,
nicht bloß als ein negatives Paley-Wiener Numerikresultat.

\subsection{Vorgeschlagene Rigoroshierarchie}
\label{sec:rigor}

Der im Baum-Meta-Analyse-Programm (tree meta-analysis) propagierten Vier-Klassen-Hierarchie
(Theorem-Ebene / bedingt / diagnostisch / spekulativ) folgend, kann jede
Komponente des Branch-Lokalitäts-Programms wie folgt klassifiziert werden:

\begin{center}
\begin{tabular}{@{}lll@{}}
\toprule
Komponente & Klasse & Quelle \\
\midrule
Definition~\ref{def:hp-bl} & Theorem-Ebene & dieses Paper \\
Theorem~\ref{thm:NE-A} & Theorem-Ebene & \cite{Geiger2026RHII} \\
Theorem~\ref{thm:NE-B} ($N\leq 15$) & Theorem-Ebene & \cite{Geiger2026RHII} \\
Vermutung~\ref{conj:NE-B-full} & bedingt (conditional) & \cite{Geiger2026RHII} \\
Theorem~\ref{thm:selberg-hp} & Theorem-Ebene & \cite{Selberg1956} \\
Prop.~\ref{prop:weyl-obstruction}--\ref{prop:fredholm-orientation} & diagnostisch & \cite{Geiger2026RHIII}, \cite{MobiusGradedPrimeHub} \\
Theorem~\ref{thm:boundary} & bedingtes Nicht-Konstanz-Theorem${}^*$ & Synthese \\
Vermutung~\ref{conj:sge} (SGE) & bedingt (empirisch) & dieses Paper, \cite{SGENote} \\
\S\ref{sec:blockchain} (Blockchain-Lesart) & spekulative Brücke & dieses Paper \\
\bottomrule
\end{tabular}
\end{center}
${}^*$ bedingt nach Vermutung~\ref{conj:NE-B-full}.

\subsection{Zukünftige Richtungen}

\begin{enumerate}
\item \textbf{Formaler Beweis von Vermutung~\ref{conj:NE-B-full}.} Ein
  SVD-Dichte-Argument für $N = \infty$ ist der mit Abstand wertvollste
  nächste Schritt: Er würde Theorem~\ref{thm:boundary} von
  bedingt (conditional) auf unbedingt (unconditional) anheben.
\item \textbf{Dedekind-Zeta-Evidenz erweitern.} Der endliche
  $\mathbb{Q}(\sqrt{-5})$-Stresstest stützt die SGE-Prognose, ersetzt
  aber kein $N\to\infty$-SVD-Dichte-Argument. Der nächste Schritt ist,
  den Test über dieses einzelne endliche Zertifikat hinaus auf weitere
  Zahlkörper und in Richtung einer echten unendlich-dimensionalen
  Grenzaussage zu erweitern.
\item \textbf{Ihara-Zeta und Yang-Lee Wurzeln.} Iharas Fall ist das
  graphentheoretische Analogon zu Selberg und es wird durch
  Vermutung~\ref{conj:sge} erwartet, dass er $\HBL = \mathrm{YES}$ erfüllt, wobei
  die Bass-Hashimoto Matrix als $H_X$ dient. Yang-Lee-Nullstellen sind demgegenüber
  durch die Temperatur parametrisierte Partitionsfunktions-Wurzeln und sollten,
  nach derzeitiger Evidenz, in der Flow-Zeta-Klasse von
  Abschnitt~\ref{sec:three-class-taxonomy} angesiedelt sein, als ein
  statistisch-mechanischer Verwandter des CRM.

\item \textbf{Bevölkerung des Zoos.} Die sechs bisher präsentierten Zoobewohner
  (Riemann, Selberg, Prime-Hub, CRM, Dedekind
  $\mathbb{Q}(\sqrt{-5})$, Ihara-Petersen) plus das Begleitpaper zum Dirichlet-Atlas
  \cite{DirichletAtlas} sind illustrativ, nicht erschöpfend.
  Ein Scan der umgebenden Literatur und der Begleit-Baumdokumente
  \cite{CandidatesMasterList} identifiziert acht weitere
  Familien, die sauber in die Drei-Klassen-Taxonomie passen und die
  SGE-Dichotomie schärfen:
  \begin{itemize}
  \item \textbf{Automorphe $L(\pi, s)$ für $\mathrm{GL}_n$}
  (Langlands-Programm): geometrisch-kompakt, \textbf{Kandidaten-SGE-YES}.
  Transfer-Algebra ist die Hecke-Algebra; Casimir- und
  Spurformel-Provenienz der umgebenden reduktiven Gruppe machen dies zu einer
  natürlichen Kandidatenklasse auf der YES-Seite
  \cite{Langlands1970,CogdellPS1994}. Das ist noch kein bereits
  gesicherter Hilbert-P\'olya-Operator für Nullstellenordinaten.
  \item \textbf{Hecke $L(s, \lambda)$ über Zahlkörpern}:
  arithmetisch, \textbf{SGE-NO}. Schwesterspezies zu Dirichlet und
  Dedekind --- Primideal-Halbgruppe plus Charakter-Twist
  \cite{IwaniecKowalski}.
  \item \textbf{Artinsche $L$-Funktionen}: geometrisch-kompakt,
  \textbf{SGE-YES}. Transfer-Algebra ist eine endliche Galoisgruppe;
  Prototyp der Galois-getriebenen YES-Seite
  \cite{Artin1931,Chebotarev1926}.
  \item \textbf{Ruelle-Anosov-Zeta $\zeta_\varphi(s)$} für hyperbolische
  Flüsse: fluss-nichtkompakt, \textbf{SGE-YES}. Direkte Generalisierung
  von CRM über den $\PSL_2(\mathbb{R})$-Boost Fall hinaus; siehe
  \cite{DyatlovZworski2016,DyatlovZworski2019}.
  \item \textbf{Selberg auf dem nicht-kompakten Quotienten}
  $\mathrm{PSL}_2(\mathbb{R})/\mathrm{SO}(2) = \mathbb{H}$:
  fluss-nichtkompakt, \textbf{SGE-YES}. Kontinuierliches Spektrum plus
  Eisensteinreihen. Algebraisch fast identisch zu CRM,
  dynamisch jedoch distinkt.
  \item \textbf{Spektren der Quasinormalmoden schwarzer Löcher}
  (\cite{DeClerckHartnollYang2025}): fluss-nichtkompakt,
  \textbf{SGE-YES}. Streuresonanzen auf asymptotisch
  de~Sitter-Hintergründen; die Wheeler--DeWitt-Wellenfunktionen der 5d
  BKL-Dynamik weisen Strukturen automorpher $L$-Funktionen und
  komplexer Primon-Gase (complex primon gases) auf, die zeta-ähnliche Spektren mit Quantengravitations-Regimes verknüpfen.
  Die erste Flow-Zeta-Funktion mit realistischer Hoffnung auf experimentellen Zugang
  via Beobachtungen des Ringdown von Gravitationswellen (LIGO/Virgo/Einstein-Teleskop).
  \item \textbf{Yang-Lee-Nullstellen}:\\[-0.6ex]
  ferromagnetische Partitionsfunktionen;
  konjektural fluss-nichtkompakt, \textbf{SGE-YES}. Temperatur als
  kontinuierlicher Parameter; das Lee-Yang Kreistheorem spielt die Rolle
  eines RH-Analogons. Statistisch-mechanischer Cousin von CRM.
  \item \textbf{Bost-Connes System} (1995): eine quantenstatistische
  Realisierung von $\zeta(s)$ als Partitionsfunktion mit spontaner
  Symmetriebrechung bei $\beta = 1$. Sitzt auf der Brücke zwischen den
  arithmetischen und Fluss-Klassen und verbindet den Zoo direkt mit
  Connes' adelischem Framework \cite{BostConnes1995,Connes1999, Connes2026}.
  \end{itemize}
  Jede dieser Familien verdient beizeiten ein eigenes FST-Mathematics Supplement.
  Eine zweite Ebene von zwölf weiteren Konstruktionen
  (Epstein, Witten-Zeta, Multiple-Zeta-Werte, Hurwitz, Lerch,
  Arakelov, motivische $L$-Funktionen, chirale Zeta, RG-Zetas, Gromov-Witten,
  Reidemeister-Torsion, QFT spektrale Zetas) und eine separate Spalte
  biologischer Muster-A-abgeleiteter Zetas (Kauffman-Attraktor,
  Proteinfaltungs-Spektren, Replikator-Jacobi-Spektren) sind in
  dem Begleit-Baumdokument katalogisiert \cite{CandidatesMasterList}.

\item \textbf{Biologische Pattern-A-Instanzen als Brücken.} Das
  FST-Biology Supplement-Programm \cite{FSTBiology} entwickelt
  Muster A in molekularen, chemischen und verhaltensbiologischen Domänen
  (Proteinfaltung als Nash-Gleichgewicht, Genregulationsnetzwerk-Attraktoren,
  Co-Evolution des Immunsystems). Diese sind vornehmlich
  Bewohner der physikalischen Seite des Programms, erzeugen jedoch
  eigene zeta-ähnliche Zählfunktionen (Kauffman-Attraktor-Zetas für Boolesche Gen-Netzwerke, Proteinfaltungs-Spektral-Zetas,
  Replikator-Jacobi-Spektren), die prinzipiell eigene
  Gehege im Zoo besiedeln könnten. Ihre primäre Klassifikation
  findet auf der physikalischen Seite von FST statt; ihre mathematische Struktur ist
  in der Begleitliste der Kandidaten (candidate document) aufgeführt.

\item \textbf{Adelisches Upgrade: Riemann revidiert.} Connes' adelisches
  Framework \cite{Connes1999, Connes2026} ersetzt die klassische
  Primzahlpotenz-Halbgruppe $\{p^m\}$ durch die Ideleklassengruppe
  $\mathbb{A}_K/K^\times$, welche eine Gruppe \emph{ist}. Falls dieses
  Programm abschließt, migriert Riemann von der SGE-NO Spalte
  (klassische Primverschiebung-Transfer-Algebra) zur SGE-YES Spalte
  (adelische Transfer-Algebra), ohne Theorem~\ref{thm:boundary} zu widersprechen ---
  die SGE-Vermutung beschreibt
  präzise, was ein erfolgreiches Hilbert-P\'olya Programm
  erreichen muss.

\item \textbf{Explizite Prime-Hub Konstruktion.} Jede künftige
  Konstruktion muss explizit durch die
  Propositionen~\ref{prop:weyl-obstruction}--\ref{prop:fredholm-orientation} navigieren.
  Unter der Vermutung~\ref{conj:sge} gelingt der Konstruktion von
  $G_\mathrm{prime}$ die Bereitstellung eines Hilbert-P\'olya Operators
  genau dann, wenn die resultierende Zyklus-Verkettungs-Algebra eine Gruppe ist.
  Zusätzlich müssen OP5.3 und OP5.5 in einen Euler-/Möbius-Kanal und
  einen globalen Spektralkanal getrennt oder durch eine echte
  Superdeterminante codiert werden; eine nackte positive
  Fredholm-Determinante hat an Zeta-Nullstellen die falsche Orientierung.

\item \textbf{Boundary-Taxonomie-Ledger.} Neuere Transferchecks machen
  aus dem Zoo nicht nur eine Liste von Spezies, sondern ein
  Guardrail-Ledger. Vier Boundary-Typen sollten für künftige Einträge
  mitgeführt werden:
  \begin{center}
  \scriptsize
  \begin{tabular}{@{}p{0.20\textwidth}p{0.32\textwidth}p{0.32\textwidth}@{}}
  \toprule
  Typ & Test & Guardrail \\
  \midrule
  Prunable appendage &
  Ihara-artige Seitenkorridore vor der Klassifikation auf ihren Core prunen &
  Endliche Baum-Anhängsel, die Zeta-/Operatorobjekte unverändert lassen, sind keine neue Spezies \\
  Schur-Cofactor obstruction &
  Für $G=(\begin{smallmatrix}A&b\\ b^t&c\end{smallmatrix})$ vollen Rang von $A$ fordern und $s=c-b^tA^{-1}b$, $Q_B=\beta^2/s$ prüfen &
  $\det(G)\neq0$ ist nur bei vollrangigem Source-Block und nichtverschwindender Kopplung ein Kurztest \\
  Fixed target vs global transfer &
  Feste Zielwert-Proben von wachsenden Fenstern und Grenzaussagen trennen &
  Das Treffen ausgewählter Ziele ist noch kein globaler Hilbert-P\'olya- oder SGE-Transfer \\
  Local pass vs global section &
  Lokale Root-, Branch- oder Verifier-Passes zu einer kohärenten niedrigen globalen Sektion zusammensetzen &
  Lokale Passes ohne Sektion bleiben diagnostisch \\
  \bottomrule
  \end{tabular}
  \end{center}
  Vorgeschlagene Prüffelder sind
  \texttt{survives\_core\_pruning}, \texttt{schur\_scalar},
  \texttt{fixed\_target\_probe}, \texttt{growth\_window\_status},
  \texttt{global\_limit\_status} und \texttt{selection\_rule}.
  Dieses Ledger ändert Theorem~\ref{thm:boundary} nicht; es verhindert,
  dass künftige Zoo-Populationen lokale oder endliche Zertifikate als
  globale Transferaussagen überlesen.

\item \textbf{Unabhängige Replikation und ein Gutachter-Verifizierbarkeitsindex (reviewer-verifiability index).}
  Die Trilogie \cite{Geiger2026RHII,Geiger2026RHIII}
  umfasst derzeit ca. 82 PDF-Seiten, 76 Theorem-Objekte, eine achtschrittige Reduktionskette, 33 CAP-Zertifikate und
  zwanzig Python-Skripte. Das Critique-Review
  \cite{RHCritiqueReview2026} identifiziert, obgleich es keine fatale Lücke findet,
  dokumentations-ergonomische Schwächen (K6): Fehlen eines
  zentralen Symbol-Glossars, kein Abhängigkeits-Graph (DAG) der Schritte A1--A8, keine vereinheitlichte
  \texttt{reproduce.sh} Pipeline, kein „reviewer fast path“ Hinweis.
  Wir schlagen einen begleitenden \emph{Reviewer Kit} Zenodo-Record vor,
  der (a) ein Symbol-Glossar, (b) einen Abhängigkeits-DAG, der
  Lemmas mit Skripten und A-Schritten verknüpft, (c) eine Datei
  \texttt{scripts/REPRODUCE.md} mit Aufrufen als Einzeiler und
  erwarteten Laufzeiten, (d) ein Lean4-Gerüst für das Shift-Parity-Lemma ---
  dessen rein algebraische Struktur ($N \geq 3$, geschlossene Form für Determinante und Spur) es zum natürlichen ersten Ziel
  für eine ernsthafte Formalisierung macht --- und (e) einen zweiseitigen „reviewer fast path“ umfasst, der darlegt, wie jeder A-Schritt innerhalb eines
  begrenzten Zeitbudgets unabhängig verifiziert werden kann. Ein beigefügter \emph{Gutachter-Verifizierbarkeits-Index} pro A-Schritt (analytisch / CAP / hybrid, geschätzte Personenstunden
  für unabhängige Verifikation) würde die Replikationskosten der Trilogie
  explizit machen. Entscheidend ist, dass dieses Programm keine
  Änderung an der Trilogie selbst erfordert; es ist eine additive Ergänzung, die
  K6 adressiert, ohne den Beweis anzutasten.
\end{enumerate}

% ============================================================
\section*{KI-Offenlegung}

Dieses Manuskript wurde mit umfangreicher KI-gestützter Unterstützung
bei Literaturorganisation, Entwurf, Übersetzung, Konsistenzprüfung,
LaTeX-Pflege und mathematischen Review-Prompts erstellt. KI-Systeme
wurden als Werkzeuge eingesetzt und sind keine Autorinnen oder Autoren,
keine Koautorinnen oder Koautoren, keine Danksagungsempfänger, keine
Wahrheitsinstanzen und keine unabhängigen Validierungsinstanzen. Alle
mathematischen Aussagen, Interpretationen und Veröffentlichungsentscheidungen
liegen ausschließlich in der Verantwortung des Autors.

% ============================================================
\begin{thebibliography}{99}

% --- Interne Projektreferenzen ---

\bibitem{Geiger2026RHI}
L.~Geiger, \emph{The Riemann Hypothesis via the v2.0 Method,
Part I: Foundations and Obstructions}, Preprint (2026), Zenodo-Konzept-DOI
\href{https://doi.org/10.5281/zenodo.19035640}{10.5281/zenodo.19035640}.

\bibitem{MetaGalerkinBridge}
L.~Geiger, \emph{A Meta Galerkin Bridge: Diagnostic Decomposition
of RH-Type Programmes}, Arbeitsentwurf (2026),
\texttt{paper/META\_GALERKIN\_BRIDGE.tex} im archivierten Vorgänger-Baum. Appendix I sammelt
die zehn essayistischen Motivationen, auf die hier verwiesen wird, inklusive der
Themen I8--I10 zur Blockchain.

\bibitem{RFEPFramework}
L.~Geiger, \emph{Functional Stability Theory --- Physical
Instantiation via the Renormalized Free Energy Principle},
Preprint, Begleitpaper (2026), Zenodo-Konzept-DOI
\href{https://doi.org/10.5281/zenodo.19036190}{10.5281/zenodo.19036190};
öffentliches Repository
\href{https://github.com/research-line/functional-stability-theory/tree/main/masters/spectrum-duality}{research-line/functional-stability-theory/masters/spectrum-duality}.
Die aktuelle Live-Version v1.7 steht noch unter dem RFEP-/Spectrum-Duality-Titel;
das v1.8-Terminologie-Update ist offen.

\bibitem{FSTYangMills}
L.~Geiger, \emph{FST-Physics: Yang--Mills Mass Gap via Pattern A},
Preprint (2025/2026), Zenodo DOI
\href{https://doi.org/10.5281/zenodo.19087433}{10.5281/zenodo.19087433}.

\bibitem{FSTTurbulence}
L.~Geiger, \emph{FST-Physics: Kolmogorov K41 Spectrum via
Pattern A}, Preprint (2025/2026), Zenodo DOI
\href{https://doi.org/10.5281/zenodo.19056813}{10.5281/zenodo.19056813}.

\bibitem{Geiger2026RHII}
L.~Geiger, \emph{A Conditional Reduction of the Riemann Hypothesis to
Even Dominance of the Weil Quadratic Form}, Preprint (2026),
Zenodo-Konzept-DOI
\href{https://doi.org/10.5281/zenodo.19035640}{10.5281/zenodo.19035640}.

\bibitem{Geiger2026RHIII}
L.~Geiger, \emph{The Riemann Hypothesis via the v2.0 Method,
Part III: Conclusio and Open Problems}, Preprint (2026),
Zenodo-Konzept-DOI
\href{https://doi.org/10.5281/zenodo.19035640}{10.5281/zenodo.19035640}.

\bibitem{DirichletAtlas}
L.~Geiger, \emph{A Weil-Kernel Numerical Atlas for Dirichlet
$L$-Function Sector Gaps: Galerkin Diagnostics and the Boundaries of
Leading-Order Theory}, Preprint (2026), Zenodo-Konzept-DOI
\href{https://doi.org/10.5281/zenodo.19960809}{10.5281/zenodo.19960809};
Quellpfad \texttt{CORE/zoo-mapping/paper/DIRICHLET\_CHARACTER\_ATLAS\_v1\_en.tex}.

\bibitem{PaleyWienerDirichlet}
L.~Geiger, \emph{Paley-Wiener Transfer for Dirichlet $L$-functions:
First-Order Perturbation and Numerical Validation}, Arbeitsnotizen
(2026),
\texttt{CORE/zoo-mapping/PALEY\_WIENER\_DIRICHLET\_TRANSFER.md}.

\bibitem{SGENote}
L.~Geiger, \emph{The Semigroup-Group Equivalence: A Unifying
Conjecture for Branch-Locality}, Arbeitsnotizen (2026),
\texttt{paper\_NE\_B\_BOUNDARY/SEMIGROUP\_GROUP\_EQUIVALENCE.md}.

\bibitem{DedekindNEBTest}
L.~Geiger, \emph{Dedekind NE-B Analog Test for
$\mathbb{Q}(\sqrt{-5})$: First Empirical Probe of the
Semigroup-Group Equivalence}, Computernotizen (2026),
\texttt{paper\_NE\_B\_BOUNDARY/\_results/DEDEKIND\_NE\_B\_TEST.md}.

\bibitem{IharaPetersenTest}
L.~Geiger, \emph{Ihara-Zeta SGE YES-side Test for the Petersen
Graph: Bass-Hashimoto, Ramanujan, and the Laplacian as HP-Operator},
Computernotizen (2026),
\texttt{paper\_NE\_B\_BOUNDARY/\_results/IHARA\_PETERSEN\_SGE.md}.

\bibitem{MobiusGradedPrimeHub}
L.~Geiger, \emph{Möbius-Graded Prime-Hub: Determinant Orientation,
Exterior-Fock Normal Form, and the OP5 Two-Channel Split}, Computernotiz
(2026),
\texttt{CORE/zetazoo/\_results/MOBIUS\_GRADED\_PRIMEHUB\_TOY.md}.

\bibitem{PrimeWeilBridgeDefect}
L.~Geiger, \emph{Prime--Weil Bridge Defect: Diagnostic and Null-Free
Fourier-CCM Operator Controls}, Computationsartefakt (2026),
\texttt{CORE/zetazoo/\_results/PRIME\_WEIL\_BRIDGE\_DEFECT.md} und
\texttt{CORE/zetazoo/\_results/PRIME\_WEIL\_BRIDGE\_DEFECT\_OPERATOR\_FOURIER.md}.

\bibitem{CRMFlowZetaNote}
L.~Geiger, \emph{CRM as Flow-Zeta: Reformulation of the Cooperative
Renormalization Model of Dark Energy as a Ruelle-Dyatlov-Zworski
Zeta on the Hyperbolic Plane}, Arbeitsnotiz (2026),
\texttt{04\_FST\_COSMOLOGY/CRM\_AS\_FLOW\_ZETA.md}.

\bibitem{Zookeeper2026}
L.~Geiger, \emph{The Spectral Zookeeper: The Riemann Hypothesis via
Spectral Microcluster Closure in the Connes--Consani--Moscovici
Framework}, Preprint (2026), Zenodo DOI
\href{https://doi.org/10.5281/zenodo.19673126}{10.5281/zenodo.19673126}.

\bibitem{CRMV}
L.~Geiger, \emph{Cooperative Renormalization Model of Dark Energy V:
The Saturation Theorem}, Preprint (2026), Zenodo-Konzept-DOI
\href{https://doi.org/10.5281/zenodo.18728935}{10.5281/zenodo.18728935}.

\bibitem{DyatlovZworski2019}
S.~Dyatlov und M.~Zworski, \emph{Mathematical Theory of Scattering
Resonances}, Graduate Studies in Mathematics \textbf{200}, American
Mathematical Society, Providence, RI (2019),
DOI~\href{https://doi.org/10.1090/gsm/200}{10.1090/gsm/200}.

\bibitem{CandidatesMasterList}
L.~Geiger, \emph{Zeta Zoo Candidates --- Master List}, Begleitendes
Baumdokument (2026),
\texttt{CORE/CANDIDATES\_MASTER\_LIST.md} im
\texttt{.ZETA-ZOO} Projektbaum. Führt den vollständigen Katalog der Zeta-Typ-Familien-Kandidaten,
identifiziert über RH v2.1, EFP-Taxonomy, V2-Physics-Testcases, FST-Biology und
ergänzende Ideenanhang-Quellen.

\bibitem{FSTBiology}
L.~Geiger, \emph{Functional Stability Theory --- Biological
Pattern-A Instances (FST-I Thermodynamic, FST-II Chemical, FST-III
Biological Stability)}, Preprint-Reihe (2026), Zenodo-Konzept-DOIs
\href{https://doi.org/10.5281/zenodo.20130544}{10.5281/zenodo.20130544},
\href{https://doi.org/10.5281/zenodo.20130563}{10.5281/zenodo.20130563}
und \href{https://doi.org/10.5281/zenodo.20130573}{10.5281/zenodo.20130573}.
Dreiteilige Publikationsreihe über skaleninvariante Muster-A-Realisierungen (Pattern-A)
von der molekularen bis zur verhaltensbasierten Biologie.

% --- Externe klassische Referenzen (verifiziert via WebSearch 2026-04-16) ---

\bibitem{Polya1982letter}
G.~P\'olya, Brief an A.~M.~Odlyzko, 3. Januar 1982, in dem eine
Göttinger Konversation mit E.~Landau (ca.\ 1912--1914) über eine
mögliche operatorentheoretische Erklärung für die Riemannsche Vermutung
geschildert wird. Bewahrt von Odlyzko auf
\url{https://www-users.cse.umn.edu/~odlyzko/polya/}.

\bibitem{Montgomery1973}
H.~L.~Montgomery, \emph{The pair correlation of zeros of the zeta
function}, in: \emph{Analytic Number Theory} (Proc.\ Sympos.\ Pure
Math.\ \textbf{XXIV}, St.\ Louis 1972), S.~181--193, American
Mathematical Society, Providence, RI (1973). Früheste publizierte
Erwähnung der Hilbert--P\'olya-Anregung als Folklore.

\bibitem{Berry1985}
M.~V.~Berry, \emph{Semiclassical theory of spectral rigidity},
Proc.\ Roy.\ Soc.\ London A \textbf{400}, 229--251 (1985), DOI
\href{https://doi.org/10.1098/rspa.1985.0078}{10.1098/rspa.1985.0078}.

\bibitem{BerryKeating1999}
M.~V.~Berry und J.~P.~Keating, \emph{The Riemann zeros and
eigenvalue asymptotics}, SIAM Rev.\ \textbf{41} (2), 236--266 (1999),
DOI
\href{https://doi.org/10.1137/S0036144598347497}{10.1137/S0036144598347497}.

\bibitem{Connes1999}
A.~Connes, \emph{Trace formula in noncommutative geometry and the
zeros of the Riemann zeta function}, Selecta Math.\ (New Series)
\textbf{5} (1), 29--106 (1999), DOI
\href{https://doi.org/10.1007/s000290050042}{10.1007/s000290050042}.

\bibitem{Connes2026}
A.~Connes, \emph{The Riemann Hypothesis: Past, Present and a
Letter Through Time}, arXiv Preprint (2026), arXiv:2602.04022
[math.NT], \url{https://arxiv.org/abs/2602.04022}.

\bibitem{ConnesConsaniMoscovici2025}
A.~Connes, C.~Consani, und H.~Moscovici, \emph{Zeta Spectral
Triples}, arXiv Preprint (2025), arXiv:2511.22755 [math.NT],
\url{https://arxiv.org/abs/2511.22755}. Begleitpaper zu
\cite{Connes2026}.

\bibitem{GeigerObstructionClassifiers}
L.~Geiger, \emph{Obstruction Classifiers for RH-Type Problems},
Arbeitspaper (2026), vgl.\ \texttt{\_archive/old\_obstruction\_classifier\_wrapper/Obstruction\_Classifiers\_v1\_en.tex}
im archivierten Vorgängerbaum.

\bibitem{SGEControlExperiment2026}
L.~Geiger, \emph{SGE Control Experiment --- Discriminating Test},
Arbeitsnotiz und Rechenartefakt (2026),
\texttt{CORE/zetazoo/\_results/SGE\_CONTROL\_EXPERIMENT.md},
erzeugt durch
\texttt{CORE/zetazoo/\_scripts/sge\_control\_experiment.py}.
Server-zertifizierter Durchlauf (ellmos-services, 2026-04-17), der die
Skalierung der Zentralisator-Dimension für zyklische $\mathbb{Z}/N$,
elementar-abelsche $(\mathbb{Z}/2)^k$ und zufällig-symmetrische
Baselines dokumentiert.

\bibitem{RHCritiqueReview2026}
L.~Geiger, \emph{RH v2.1 --- Systematic Critique Review (K1--K6)},
Arbeitsnotiz (2026), \texttt{SESSIONS/\_reviews/RH\_CRITIQUE\_REVIEW\_2026-04-17.md}
im \texttt{.ZETA-ZOO} Projektbaum.
Paralleles Assessment der Foundation--Even-Dominance--Conclusio
Trilogie durch vier Agenten gegenüber sechs externen Kritikpunkten; Fazit: alle teilweise berechtigt,
keine fatal.

\bibitem{EkelandTemam1999}
I.~Ekeland und R.~T\'emam, \emph{Convex Analysis and Variational
Problems}, SIAM Classics in Applied Mathematics \textbf{28}, Society
for Industrial and Applied Mathematics, Philadelphia (1999),
ISBN 978-0-89871-450-0. Standardreferenz für das Lemma von Mazur und
das Eindeutigkeitsprinzip (strikte Konvexität plus Koerzivität), das in
Lemma~\ref{lem:ucu} verwendet wird.

% --- Hinzugefügte externe klassische / kontemporäre Referenzen (WebSearch-verifiziert 2026-04-17) ---

\bibitem{KeatingSnaith2000}
J.~P.~Keating und N.~C.~Snaith, \emph{Random Matrix Theory and
$\zeta(1/2 + it)$}, Commun.\ Math.\ Phys.\ \textbf{214} (2000),
S.~57--89, DOI~\href{https://doi.org/10.1007/s002200000261}{10.1007/s002200000261}.

\bibitem{Bombieri2000}
E.~Bombieri, \emph{Problems of the Millennium: The Riemann
Hypothesis}, Offizielle Problembeschreibung des Clay Mathematics Institute (2000),
\url{https://www.claymath.org/wp-content/uploads/2022/05/riemann.pdf}.

\bibitem{Conrey2003}
J.~B.~Conrey, \emph{The Riemann Hypothesis}, Notices of the AMS
\textbf{50} (2003), no.~3, S.~341--353,
\url{https://www.ams.org/notices/200303/fea-conrey-web.pdf}.

\bibitem{BostConnes1995}
J.-B.~Bost und A.~Connes, \emph{Hecke algebras, type III factors
and phase transitions with spontaneous symmetry breaking in number
theory}, Selecta Math.\ (New Series) \textbf{1} (1995), no.~3,
S.~411--457,
DOI~\href{https://doi.org/10.1007/BF01589495}{10.1007/BF01589495}.

\bibitem{Langlands1970}
R.~P.~Langlands, \emph{Problems in the theory of automorphic
forms}, in: C.~T.~Taam (Hrsg.), \emph{Lectures in Modern Analysis
and Applications III}, Lecture Notes in Mathematics \textbf{170},
Springer, Berlin--Heidelberg (1970), S.~18--61,
DOI~\href{https://doi.org/10.1007/BFb0079065}{10.1007/BFb0079065}.

\bibitem{DeClerckHartnollYang2025}
M.~De~Clerck, S.~A.~Hartnoll, und M.~Yang,
\emph{Wheeler--DeWitt wavefunctions for 5d BKL dynamics, automorphic
$L$-functions and complex primon gases}, arXiv Preprint (2025),
arXiv:2507.08788 [hep-th],
\url{https://arxiv.org/abs/2507.08788}.

\bibitem{Artin1931}
E.~Artin, \emph{Zur Theorie der L-Reihen mit allgemeinen
Gruppencharakteren}, Abh.\ Math.\ Sem.\ Univ.\ Hamburg \textbf{8}
(1931), S.~292--306,
DOI~\href{https://doi.org/10.1007/BF02941010}{10.1007/BF02941010}.

\bibitem{Chebotarev1926}
N.~Tschebotareff, \emph{Die Bestimmung der Dichtigkeit einer Menge
von Primzahlen, welche zu einer gegebenen Substitutionsklasse
geh\"oren}, Math.\ Ann.\ \textbf{95} (1926), S.~191--228,
DOI~\href{https://doi.org/10.1007/BF01206606}{10.1007/BF01206606}.

\bibitem{CogdellPS1994}
J.~W.~Cogdell und I.~I.~Piatetski-Shapiro, \emph{Converse theorems
for $GL_n$}, Publ.\ Math.\ Inst.\ Hautes \'Etudes Sci.\ \textbf{79}
(1994), S.~157--214,
\url{https://www.numdam.org/item/PMIHES_1994__79__157_0/}.

\bibitem{DyatlovZworski2016}
S.~Dyatlov und M.~Zworski, \emph{Dynamical zeta functions for
Anosov flows via microlocal analysis}, Ann.\ Sci.\ \'Ec.\ Norm.\
Sup\'er.\ (4) \textbf{49} (2016), no.~3, S.~543--577,
DOI~\href{https://doi.org/10.24033/asens.2290}{10.24033/asens.2290},
arXiv:1306.4203.

\bibitem{Selberg1956}
A.~Selberg, \emph{Harmonic analysis and discontinuous groups in
weakly symmetric Riemannian spaces with applications to Dirichlet
series}, J.\ Indian Math.\ Soc.\ \textbf{20}, 47--87 (1956).

\bibitem{Weil1952}
A.~Weil, \emph{Sur les ``formules explicites'' de la th\'eorie des
nombres premiers}, Meddelanden fr\aa n Lunds Universitets
Matematiska Seminarium, Supplementband tillägnet Marcel Riesz,
S.~252--265, C.~W.~K.\ Gleerup, Lund (1952).

\bibitem{IwaniecKowalski}
H.~Iwaniec und E.~Kowalski, \emph{Analytic Number Theory}, American
Mathematical Society Colloquium Publications \textbf{53}, AMS,
Providence, RI (2004), xi+615~S., ISBN 0-8218-3633-1.

\bibitem{Titchmarsh1986}
E.~C.~Titchmarsh, \emph{The Theory of the Riemann Zeta-Function},
2. Auflage, überarbeitet von D.~R.~Heath-Brown, Oxford University Press,
Oxford (1986), ISBN 0-19-853369-1.

\end{thebibliography}

\end{document}
